\documentclass[parskip=full]{article}

%%%%%%%%%%%%%%%%%%%%%%%%%%%%%%%%%%%%%%%%%%%%%%%%%
% Basic includes
%%%%%%%%%%%%%%%%%%%%%%%%%%%%%%%%%%%%%%%%%%%%%%%%%


\usepackage{arydshln}
\usepackage{array}
\usepackage{makecell}
\usepackage{tikz}
\usepackage{amsmath}
\usepackage{amsfonts}
\usepackage{verbatim}
\usepackage{mathrsfs}
\usepackage{amssymb}
\usepackage{amsthm}
\usepackage{bbm}
\usepackage{bm}
\usepackage[capposition=top]{floatrow}
%\usepackage[dvipsnames]{xcolor}
\usepackage{ulem}
\usepackage[labelfont=bf]{caption}
\usepackage{subcaption}
\usepackage{graphicx}
\usepackage{float}
\usepackage{natbib}
\usepackage{booktabs}
\usepackage{longtable}
\usepackage{multirow}
\usepackage{rotating}
\usepackage{pdflscape}
\usepackage[title]{appendix}
\usepackage{footmisc}
\usepackage{bigdelim}
\usepackage{tabularx}
\usepackage{xspace}
\newcommand{\A}{\ensuremath{\mathcal{A}}\xspace}
\newcommand{\B}{\ensuremath{\mathcal{B}}\xspace}
\newcommand\pa[1]{\ensuremath{\left(#1\right)}}
\usepackage{epigraph}
\usepackage{setspace}

\usepackage[pdftex, colorlinks=true, linkcolor=blue, citecolor=blue, urlcolor=blue, bookmarks=false]{hyperref}
%\onehalfspacing
\setstretch{1.5}

%%%%%%%%%%%%%%%%%%%%%%%%%%%%%%%%%%%%%%%%%%%%%%%%%%%
% Microtype
%%%%%%%%%%%%%%%%%%%%%%%%%%%%%%%%%%%%%%%%%%%%%%%%%%%
\usepackage[activate={true,nocompatibility},final,tracking=true,kerning=true,spacing=true,factor=1100,stretch=10,shrink=10]{microtype}
% activate={true,nocompatibility} - activate protrusion and expansion
% final - enable microtype; use ``draft'' to disable
%tracking=true, kerning=true, spacing=true - activate these techniques
% factor=1100 - add 10% to the protrusion amount (default is 1000)
% stretch=10, shrink=10 - reduce stretchability/shrinkability (default is 20/20)

%%%%%%%%%%%%%%%%%%%%%%%%%%%%%%%%%%%%%%%%%%%%%%%%%%%
% Hyperref setup
%%%%%%%%%%%%%%%%%%%%%%%%%%%%%%%%%%%%%%%%%%%%%%%%%%%
\usepackage{url}
\makeatletter
\g@addto@macro{\UrlBreaks}{\UrlOrds}
\makeatother
\usepackage{hyperref}
\hypersetup{
  colorlinks,%
  citecolor=black,%
  filecolor=black,%
  linkcolor=blue,%
  urlcolor=black
}

%%%%%%%%%%%%%%%%%%%%%%%%%%%%%%%%%%%%%%%%%%%%%%%%%%%
% Font setup
%%%%%%%%%%%%%%%%%%%%%%%%%%%%%%%%%%%%%%%%%%%%%%%%%%%
\usepackage{mathpazo} % math & rm
%\linespread{1}        % Palatino needs more leading (space between lines)
\usepackage{parskip}
\usepackage[scaled]{helvet} % ss
\usepackage{courier} % tt

\usepackage{filecontents}

\normalfont
\usepackage[T1]{fontenc}

\usepackage{titlesec}
\titleformat{\section}
{\normalfont\fontsize{12}{17}\bfseries}
{\thesection}
{1em}
{}

\titleformat{\subsection}
{\normalfont\fontsize{11}{16}\bfseries}
{\thesubsection}
{1em}
{}

%%%%%%%%%%%%%%%%%%%%%%%%%%%%%%%%%%%%%%%%%%%%%%%%%
% macros
%%%%%%%%%%%%%%%%%%%%%%%%%%%%%%%%%%%%%%%%%%%%%%%%%
\newcolumntype{C}[1]{>{\centering\let\newline\\\arraybackslash\hspace{0pt}}m{#1}}
\newcommand{\indicator}[1]{\mathbbm{1}_{\left\{ {#1} \right\} }}
%\newcommand{\resultsloc}{../code/python/output/}
\newcommand{\subfigwidth}{0.475\textwidth}
\newcommand{\subfigwidthtwo}{0.375\textwidth}
\newcommand{\figwidth}{0.6\textwidth}
\DeclareMathOperator*{\argmax}{arg\,max}

\setlength\epigraphwidth{0.9\textwidth}
\setlength\epigraphrule{0pt}
\renewcommand{\epigraphflush}{center}
\usepackage{etoolbox}
\makeatletter
\patchcmd{\epigraph}{\@epitext{#1}}{\itshape\@epitext{#1}}{}{}
\makeatother

%%%%%%%%%%%%%%%%%%%%%%%%%%%%%%%%%%%%%%%%%%%%%%%%%
% Page and paragraph settings
%%%%%%%%%%%%%%%%%%%%%%%%%%%%%%%%%%%%%%%%%%%%%%%%%
\setlength{\parskip}{0.0pc} \setlength{\parindent}{0.25in}
\setlength{\topmargin}{-3pc}
\setlength{\textheight}{9in}
\setlength{\oddsidemargin}{0pc} \setlength{\evensidemargin}{0pc}
\setlength{\textwidth}{6.5in}
\makeatletter
\@fpsep\textheight
\makeatother

\newenvironment{fignote}{\begin{quote}\footnotesize}{\end{quote}}
\def\sym#1{\ifmmode^{#1}\else\(^{#1}\)\fi}
\renewcommand{\baselinestretch}{1.5}

% Formats the notes at the bottom of tables and figures.
\newenvironment{fignote2}[1]
{
\par
\medskip
\begin{minipage}{#1}
\footnotesize
}
{
\end{minipage}
}

% Formats the notes at the bottom of tables and figures.
\newenvironment{fignote3}[1]
{
%\par
\begin{minipage}{#1}
\footnotesize
}
{
\end{minipage}
}

%%%%%%%%%%%%%%%%%%%%%%%%%%%%%%%%%%%%%%%%%%%%%%%%%
% Title page
%%%%%%%%%%%%%%%%%%%%%%%%%%%%%%%%%%%%%%%%%%%%%%%%%
%%%%%%%%%%%%%%%%%%%%%%%%%%%%%%%%%%%%%%%%%%%%%%%%%
% Title page
%%%%%%%%%%%%%%%%%%%%%%%%%%%%%%%%%%%%%%%%%%%%%%%%%

\title{\vspace{-1.8cm}Rent Control? Bad in Theory, Not Always in Practice\footnote{We extend our gratitude to the seminar attendees at the University of Wisconsin-Madison, University of Southern California, Queen's University, York University, McMaster University, University of Windsor, as well as to Joseph Steinberg, Efraim Benmelech, Ay\c{s}e \.{I}mrohoro\u{g}lu, Lu Han, Shang Zeng, Ed Glaeser, Jonathan Alschech, Nick Sander, Boaz Abramson, Pau Pujol\`as, Shouyong Shi, Amy Sun, Yahong Zhang, Ruoyu Chen, Chenggang Zhou, and John Baker for their valuable insights and comments. We owe a very special thank-you to Joseph Steinberg who generously gave us access to one of his high-powered computers for this research. The data and replication code are available on GitHub at \url{https://github.com/srotberg/rent_control_bad_in_theory_not_always_in_practice}.}}

%For those who like audio, we have an AI-generated Spotify episode explaining the paper: \url{https://spotifycreators-web.app.link/e/QVkPeRFoUOb}.

\author{Shahar Rotberg\footnote{Canada Mortgage and Housing Corporation. Corresponding author contact: \url{shahar.rotberg@gmail.com}. Website: \url{www.sites.google.com/site/shaharrotberg}. Shahar Rotberg did this work on his personal equipment and time. He exclusively owns the code that was used to solve the model in this paper.} \and Lin Zhang\footnote{Canada Mortgage and Housing Corporation. Contact: \url{lzhang@cmhc-schl.gc.ca}. The opinions expressed herein do not necessarily represent the views of CMHC.}}

\pagestyle{plain}

\begin{document}

\maketitle

\setstretch{1.5}
\begin{abstract}
\onehalfspacing

\noindent Most economists oppose rent control, arguing that it reduces welfare because it distorts rental markets. A quantitative general-equilibrium model calibrated to Toronto shows that, consistent with the Second Welfare Theorem, rent control is welfare-reducing when the government can efficiently redistribute landlord profits to low-income renters via lump-sum taxes and transfers. However, when fiscal redistribution is imperfect to empirically-realistic degrees, common types of rent control improve long-run welfare by serving as a second-best redistributive instrument. In this setting, the equity gains from transferring surplus from landlords to low-income renters outweigh the distortions caused by rent control commonly highlighted by the literature.
%words=100
\\[1.5ex]

\noindent \textbf{JEL Classification:} E17; R31; R52. 

\noindent \textbf{Keywords:} Rent control, general equilibrium, housing, tenancy control, unit control, homelessness.
\end{abstract}
\setstretch{1.5}

\section{Introduction}

Rent control is one of the most widely implemented — and hotly debated — housing policies worldwide, typically aimed at improving housing affordability for low-income renters.

Economists have long opposed rent control, arguing that it discourages rental housing supply, misallocates units, and generates deadweight loss \citep{gl-2003}. In standard models, this conclusion follows under either of two conditions: perfectly elastic long-run rental supply or efficient lump-sum redistribution. In particular:

\begin{itemize}
\item If long-run rental supply is perfectly elastic, rent control must reduce long-run welfare. Any binding ceiling that drives rents below the break-even level causes rental supply to collapse in the long run, leaving households who are unable to purchase housing without access to rental housing.
\item Even if long-run rental supply is upward sloping, rent control has no welfare-improving rationale when redistribution is efficient. In this case, the government can deregulate the rental market to maximize surplus and use lump-sum taxes and transfers to redistribute landlord profits, thereby achieving the welfare-maximizing allocation without distorting housing supply.
\end{itemize}

Both conditions are empirically challenged. Long-run rental supply is far from perfectly elastic, especially in dense urban environments where rent control is most prevalent \citep{orr-1968,cy-1994,susin-2002,fack-2006,sy-2020,rs-2023}. Moreover, redistribution through the tax-and-transfer system is often incomplete or inefficient: rental profits accrue disproportionately to wealthy households, and governments rarely tax these profits fully or transfer the revenues directly to low-income renters.

Using a general-equilibrium model with a realistically-calibrated upward-sloping rental supply, we show that, in accordance with the Second Welfare Theorem, rent control is welfare-reducing under efficient redistribution. However, when fiscal redistribution is imperfect to realistic degrees, two of the most common types of rent control — namely tenancy control and unit control — can be welfare improving: the redistributive gains to low-income households outweigh the distortions typically emphasized in the rent-control literature.

Households in the model face uninsurable income risk, are heterogeneous in labor productivity, and experience deterministic life-cycle changes in family size. They choose among the following: renting in the controlled or free market; owning a home subject to transaction costs, with long-term mortgages available as a financing option; or remaining homeless when available rental options are unaffordable or would require excessively constrained consumption.

Rent control in our model introduces several distortions that operate through household housing choices as well as via rental supply. First, rent discounts in the controlled market, together with fixed moving costs faced by all renters, generate lock-in effects. These effects increase incentives for controlled renters to remain in place as income or household size change, leading to delayed adjustment and housing misallocation.

Second, rent control restricts access to higher-quality rental units, reducing housing flexibility for renters who opt for the controlled market. When combined with lock-in effects, the inability to choose higher-quality housing in anticipation of higher future incomes or larger household size leads to more severe and persistent housing misallocation. By contrast, free-market rentals are available at both low and high quality.

Third, controlled rentals impose an income-scaled cost on tenants, capturing in a reduced-form way time-use distortions — such as constrained location choices in the regulated market. As a result, controlled tenants must absorb costs that reduce the true redistributive benefits of rent control.

Rent control can also distort rental supply, which is provided by a competitive rental company. This company faces convex marginal costs — calibrated to match empirical long-run rental supply elasticities — and chooses how many housing units to operate as rentals, taking prices and regulations as given. Only a subset of the housing stock — typically newer units — is exempt from rent control, and even within this exempt stock only a limited share can be rented, with the remainder absorbed by owner-occupiers. The firm supplies rentals from the exempt stock up to its rental utilization limit whenever free-market rents exceed marginal costs. If exemptions are sufficiently large, this capacity suffices to achieve the efficient rental quantity. When exemptions are limited, however, the firm supplies additional units in the controlled segment as long as the marginal cost of supplying a controlled unit is weakly below the expected controlled rent. As a result, deeper rent discounts reduce the rental supply by tightening the firm’s profitability condition unless exemptions are large enough to offset this effect. To match empirical landlordship patterns from the 2019 Survey of Financial Security, rental profits accrue to the top 35\% of the wealth distribution. 

To test whether our welfare-improving result holds across policy designs, we examine the two most common rent‑control regimes — tenancy control and unit control — which differ in their allocation mechanisms and eviction implications.

Under tenancy control, landlords can raise rents freely upon tenant turnover, allowing rents paid by new tenants to clear the market — a mechanism that prevents shortages. Since households are free to choose between the controlled and free rental markets, competition for controlled units results in an "entry premium" relative to free-market rents that new tenants must pay to access these units. Rent discounts, therefore, only accumulate gradually as tenants remain longer in the unit. Because tenancy control laws explicitly permit rent increases upon turnover, this "entry premium" is not a modeling artifact but a legally and empirically relevant market-clearing object \citep{arnott-1995}, which we document in our 2016 Toronto data. 

Under unit-based control, the regulated rent is fixed at a constant level below the free-market rate for the duration of a tenancy, regardless of tenant turnover. This provides immediate discounts to new tenants but generates excess demand and shortages. Allocation therefore occurs through non-price mechanisms such as lotteries, although households always retain the option to instead rent in the free market, buy a home or remain homeless, a modeling approach used elsewhere in the literature \citep{fmn-2023}. 

These setups imply three key tradeoffs between unit control and tenancy control. The first concerns how controlled units are allocated. Under tenancy control, prices determine who receives a controlled unit, whereas under unit control allocation occurs through lotteries. As a result, tenancy control is likely to allocate units more efficiently across new tenants. On the other hand, tenancy control may generate greater dynamic misallocation: as rent discounts grow with tenure under tenancy control, renters have stronger incentives to remain in their units even when their income and household size change, potentially increasing misallocation relative to unit control.

The second tradeoff concerns the timing of discounts: tenancy control delivers affordability gradually through longer tenure, while unit control provides immediate rent relief but may exclude households with very high willingness to pay.

The third tradeoff concerns eviction risk. Because landlords can reset rents upon tenant turnover under tenancy control but not under unit control, tenure-dependent rent discounts widen the gap between rents paid by incumbents and rents landlords could obtain from new tenants, strengthening incentives for tenant replacement. To capture this channel, we allow eviction risk in tenancy-controlled units to increase with the size of the incumbent’s rent discount.

We calibrate the stationary equilibrium of our model to Toronto, a densely populated and persistently unaffordable city where tenancy control operates alongside exemptions that applied to stock built after November 1, 1991 during our sampling period. We evaluate welfare using the ex-ante expected utility of newborns, a criterion that is independent of incumbency \citep{ralws-1971,conesa-et-al,gkkoc}, and run four experiments in which we compare the long-run counterfactual equilibria to our benchmark. 

To study how the degree of redistribution efficiency on the part of government affects the welfare consequences of rent control relative to rent decontrol, we conduct the following three experiments. In all three, we fully decontrol the rental market and turn the redistribution system’s efficiency "knob" by increasing fiscal redistribution efficiency from one experiment to the next. In the first, additional tax revenues — generated via taxation of higher landlord profits — are rebated through income tax cuts to the entire population to satisfy the government budget constraint. In the second, excess revenues are rebated in a more targeted way via lump-sum transfers only to renters and the homeless. In the third, we consider a counterfactual with a highly efficient tax-and-transfer system: all additional landlord profits are fully taxed and redistributed in a lump-sum fashion to newborns only. 

Lastly, to test whether our results also extend to unit control, we run one more experiment. In this experiment, we compare the benchmark economy under tenancy control to a unit-control regime in which the regulated rent charged throughout a tenancy is set equal to the benchmark’s average controlled rent, and the share of the housing stock that is exempt from regulation is held fixed at its benchmark value.

Our experiments yield three main findings. First, in the first two experiments with low degrees of redistribution efficiency, full rent decontrol would reduce long-run welfare and substantially increase homelessness, accompanied by modest increases in equilibrium rents and house prices. Second, when additional rental profits are fully taxed and redistributed efficiently to newborns only, as we do in our third experiment, decontrol would raise welfare substantially, though homelessness would still slightly rise because older groups would still face higher equilibrium rents. Third, even though unit control would allocate housing less efficiently than tenancy control across income groups, it would only slightly reduce welfare compared to the benchmark with tenancy control — and still be preferred to full decontrol as executed in the first two experiments. 

Economists often oppose rent control on the grounds that it generates deadweight loss. Economic arguments in favor of rent control, therefore, typically emphasize costly transitions: even if a deregulated steady state is welfare-improving, short-run adjustment costs can make decontrol undesirable when rent control is already in place \citep{fmn-2023}. In contrast, we show that with realistically imperfect redistribution and a data-consistent upward-sloping rental supply, both tenancy and unit control can dominate full decontrol in the long run, even after accounting for distortions commonly identified by the literature. This result demonstrates that the consensus view that rent control reduces welfare does not always hold under empirically realistic conditions.

\section{Related literature}\label{sec:related_lit}

Classic textbook theory — developed primarily for unit-control regimes with no exemptions — predicts that landlords will lower the rental quantity supplied below the efficient level when rent ceilings bind, leading to persistent shortages, misallocation, and deadweight loss \citep{suen-1989,gl-2003,diamond-2019}. Based on such positive analyses, economists have typically recommended removing rent control altogether \citep{lindbeck1971,buchanan_1986,friedman_stigler_1946}. Two key assumptions underpin this normative conclusion: that rental supply is infinitely elastic in the long run — so landlords earn no economic profits — and that if rental supply is upward sloping, any profits that arise after deregulation can be costlessly and efficiently redistributed through lump-sum taxes and transfers to maximize welfare.

However, a growing empirical and structural literature finds that rental supply is far from perfectly elastic, even in the long run, especially in the biggest cities where rent control normally is found \citep{orr-1968,cy-1994,susin-2002,fack-2006,sy-2020,rs-2023}. This implies that landlords often do earn profits that, in principle, could be taxed and redistributed. Yet in practice, governments face real constraints that prevent them from fully taxing these profits and targeting the proceeds to those in need. When these fiscal limitations are combined with realistic supply elasticities, the standard recommendation to abolish rent control becomes normatively ambiguous. We take a previously unexplored approach to this normative question, employing a general equilibrium model to evaluate whether rent control can raise welfare under realistic economic conditions.

The positive (descriptive and causal) analysis of rent control is well developed and has been documented by numerous studies. Recent works by \citet{diamond-2019} and \citet{dmq-2019} exploited quasi-experimental variation from San Francisco’s expansion of tenancy control to estimate how supply, rents, and property values changed. They find that limiting rent increases for sitting tenants reduced rental supply and raised market rents for new tenants, harming newcomers. Other works by \citet{sims-2007} and \citet{app-2014} find that the repeal of rent control in 1994 in Cambridge, Massachusetts raised house prices and free-market rents. Our results for full decontrol align with the two latter studies. All of these studies focus on positive analyses, however, rather than whether rent control can raise welfare in the long run when the fiscal authority is limited in its ability to redistribute resources.

Moving beyond textbook treatments, the theoretical literature first analyzed rent control in partial equilibrium and subsequently extended the analysis to general-equilibrium models. \citet{fallis-1985} offered an early static partial equilibrium analysis of Toronto’s tenancy-controlled rental market with exemptions, emphasizing the importance of partial exemptions for newer buildings. \citet{be-2000} developed a static model of tenancy control with asymmetric information, in which landlords set entry rents without knowing a tenant’s expected stay. We extend these insights by embedding tenancy control in a dynamic general-equilibrium model. The model features key wedges that generate misallocation of the rental stock and weaken the intended redistributive benefits of rent control, as well as a rental supply structure that increases with exemptions and declines with rent discounts.

In contrast to these partial equilibrium models, \citet{eric-1995} and \citet{fmn-2023} analyze regimes more akin to unit control in general equilibrium with a dynamic lifecycle,\footnote{In \citet{fmn-2023}, rent control grants larger rent discounts to longer-staying tenants, but still offers new movers immediate discounts relative to free-market rents. As a result, shortages emerge in their model, and lotteries are used to clear the controlled rental market. Therefore, their setup aligns more closely with the structure of unit control than with tenancy control, although it combines both to a degree.} but do not study tenancy control. These studies evaluate welfare along the transition towards decontrol. Our quantitative macro-housing model focuses on showing that rent control can enhance welfare even in the long run if rental supply is upward sloping and the government’s ability to redistribute excess revenues is limited.

While the market-clearing mechanisms operating under unit control regimes have long been understood, several more recent empirical studies examine how tenancy control clears the market. \citet{nagy-1997}, \citet{early-2000}, and \citet{cjq-2023} find that new tenants in regulated units often pay an "entry premium" relative to rents of comparable free-market units. We confirm this pattern in our Toronto data, implying that under tenancy control, markets clear through prices paid by new tenants rather than through rationing. While \citet{arnott-1995} has made this point theoretically in the past, our model formalizes this mechanism by endogenizing the "entry premium": new tenants willing to pay upfront receive tenure-dependent discounts over time. This creates a price-based allocation mechanism absent under unit control, which as we show improves the allocation of controlled units across the income distribution relative to unit control. On the other hand, we also show that because rent discounts under tenancy control accrue gradually, this regime worsens the allocation of discounts across the age distribution relative to unit control, which delivers immediate benefits to new tenants.

Beyond the standard economic concerns typically raised in discussions of rent control, we link our analysis to another major and pressing social issue that has often been thought to be connected to rent control: homelessness. Early analyses by \citet{tucker1987,tucker1989heritage,tucker1991cato,tucker1991public} estimated that rent control increased homelessness. However, a later study by \citet{eo-1998}, which to our knowledge remains the most recent attempt to causally estimate the effect of rent control on homelessness using reliable micro-data, found that rent control either had no effect on homelessness or reduced it. Recent macro-housing models \citep{az-2024, abramson-2024} examine the effects of housing policies — other than rent control — on homelessness. Our study quantifies how homelessness responds to various rent control policies, and our experiments that find that both tenancy and unit control reduce homelessness relative to full decontrol are consistent with the findings of \citet{eo-1998}. 

\section{Model}\label{sec:model}

\subsection{General overview}\label{sec:brief_overview}

The model features overlapping generations of finitely lived households, along with price-taking, profit-maximizing rental and housing construction firms, and federal and local governments. Households are heterogeneous with respect to age and household size, labor market productivity, housing tenure (ownership, renting, and homelessness),\footnote{Beyond our substantive interest in homelessness, its inclusion in our model is essential. Without it, some low-income households would appear so consumption-constrained after paying for housing that even small rent reductions would generate unrealistically large welfare gains, driven by extremely high marginal utility of consumption of the very poor. Imposing a minimum income to ensure all households can afford at least the smallest rental is, in our view, a less satisfactory alternative, as the arbitrary choice of that threshold would itself distort welfare comparisons and bias results.} housing assets, mortgage debt, and financial assets. Renters are further differentiated by the type of rental market in which they participate — either controlled or free — as well as by the duration of their residence in their unit. Working-age households participate in the labor market, while retirees receive government pensions that depend on their average lifetime earnings.

Controlled units are supplied only at low quality, whereas free-market rentals are available at both low and high quality. Renting in the controlled market comes with an explicit cost that scales up with labor income for working-age households, which is meant to capture time-use distortions such as longer commuting times due to worse locations of controlled units. There are two types of rent control in our model, which we examine separately in our experiments: tenancy control and unit control.

Under tenancy control, renters freely choose between the controlled and free rental markets and optimally determine their tenancy duration. However, because landlords can reset rents upon tenant turnover, competition for these units increases the initial rent paid by new tenants relative to the free-market rent. This "entry premium" clears the market for controlled units by equating aggregate demand for controlled rentals with supply. Rent discounts therefore accumulate gradually with tenure, which in turn raises the risk of eviction, capturing landlords’ stronger incentives to evict tenants when higher rents can be charged to new occupants.

In contrast, under unit control, regulated rents are fixed below free-market rates for the duration of a tenancy, regardless of tenant turnover, generating excess demand and shortages. Access to controlled units is therefore rationed through lotteries, and the market-clearing object becomes the probability of having an opportunity to rent a controlled unit. Households that win the lottery can choose to rent a controlled unit and receive an immediate rent discount. However, they are not forced to rent in the controlled market, and may instead rent in the free market, purchase a home, or remain homeless. Households that lose the lottery cannot access the controlled market in that period, but likewise may rent in the free market, buy a home, or remain homeless. Once in a unit-controlled rental, tenants can choose to remain indefinitely, though rents do not decline with tenure. Since regulated rents cannot be reset for new tenants, there is no eviction risk under unit control.

All renters face a fixed cost when moving between rental units, while homeowners pay proportional property taxes and maintenance costs, and can finance housing purchases using long-term mortgage contracts, following \citet{kkk-2021}.\footnote{Our general equilibrium results depend on how easily households can switch between renting and owning, as this affects rental demand and, consequently, equilibrium rents. Because mortgage contract features critically affect tenure decisions, we model mortgages as realistic long-term contracts.} 

A competitive, profit-maximizing rental firm supplies rental housing to both the free and controlled segments, facing convex marginal costs, calibrated to generate realistic rental supply elasticities. Rent control applies only to the regulated portion of the housing stock, while the government can exempt from regulation a subset of the stock — interpreted as newer housing.

Housing is supplied endogenously by a competitive, profit-maximizing construction company, which is subject to increasing marginal costs. In equilibrium, the housing price equates the aggregate demand for housing (from both homeowners and renters) to the aggregate supply of housing.

The model incorporates both local and federal governments. The local government offers local public services to residents and shelters to the homeless, financed through property taxes. The federal government funds pensions, lump-sum transfers, and federal public goods through income and consumption taxes.

Because the model economy represents a single city, we take the interest rate, $r$, as exogenous. 

\subsection{Demographics}\label{subsection:demographics}

Households enter the model at age $1$, retire from the labor force at age $J_{R}$, and live for $J$ periods, where age is denoted by $j$. They face no mortality risk prior to age $J$, but die with certainty afterward. Household preferences are characterized by the following flow utility function:

\begin{equation*}
    u_{j}(c,h)=\xi_{j}^{\sigma}\frac{\left(c^{1-\gamma} h^{\gamma}\right)^{1-\sigma}}{1-\sigma},
\end{equation*}

\noindent where $c$ denotes non-housing consumption and $h$ represents housing services consumption. The parameter $\gamma$ denotes the weight assigned to the consumption of housing services in the utility, $\sigma$ captures relative risk aversion and its inverse is the elasticity of intertemporal substitution, and $\xi_{j}$ is a deterministic age-specific equivalence scale that adjusts for changes in household size over their lifecycle, following \citet{lm-1980}, \citet{wz-2010}, and \citet{xu-2013}. Households discount future consumption at the rate $\beta$, and pay a proportional tax $\tau_{c}$ on their non-housing consumption.

In addition, households derive warm-glow utility from bequests they leave behind upon death following \citet{cagetti-denardi-2006}. The bequest utility function is given by $\omega_{B}(q)=\mu \frac{q^{1-\sigma}}{1-\sigma}$, where $q$ denotes net worth at the end of life (formally defined in Section \ref{subsection:household_optimization}), and $\mu$ modulates the strength of the bequest motive.

\subsection{Income and income taxes}\label{subsection:income}

Working-age households ($j<J_{R}$) supply labor inelastically. Labor income for an age-$j$ household with labor productivity components $\chi$ and $e$ is given by $y_{\ell,j}(\chi,e)=\chi e\zeta_j$, where $\chi$ remains constant over the lifecycle and is drawn from the unconditional distribution $\hat{F}(\chi)$. The idiosyncratic shock $e$ evolves according to a Markov process $\tilde{F}(e \mid e')$, and the term $\zeta_j$ captures the deterministic component of lifecycle earnings. Upon retirement, households receive pensions $y_{R}(\chi,e)$ based on their labor productivity at retirement.\footnote{Households with higher productivity at retirement are likely to have had higher earnings over their lives, and therefore receive larger pensions.}

In addition to labor income, households earn capital income on their savings: $y_{k}(a)=ar$, where $r$ is the interest rate. Taxable income is defined as:
\begin{equation*}\label{eq:taxable_labor_income}
  y^{x}_{j}(\chi,e,a) =y_{\ell,j}(\chi,e)+\nu y_{k}(a),
\end{equation*}

\noindent where $\nu$ denotes the taxable share of capital income. Taxable income is subject to a progressive tax schedule:
\begin{equation}\label{eq:income_tax}
  y^{\tau}_{j}(\chi,e,a) = y^{x}_{j}(\chi,e,a)-\tau_{\ell}y^{x}_{j}(\chi,e,a)^{\psi},
\end{equation}

\noindent where $\tau_{\ell}$ determines the average tax rate and $\psi$ governs the progressivity of the tax schedule, following \citep{hsv-2020}.

As in \citet{kkk-2021}, living households across all ages receive financial endowments each period, which are equal to the bequests from the deceased in the preceding period. These are denoted by $\vartheta (\chi,e)$, capturing a positive correlation between bequests and household income.

Households receive profits earned by rental and construction firms in a lump-sum fashion, denoted by $T(s)$, and these lump-sum transfers may depend on the household's state $s$. 

Lastly, households may also receive lump-sum transfers $T_{fed}(s)$ from the federal government, which may also vary with the household’s state $s$. We employ these transfers in counterfactual simulations in which government redistribution is more efficient than that achieved by income tax cuts to all households.

\subsection{Housing}\label{subsection:housing_choice}

Households choose between owning or renting a dwelling, denoted by $o \in \mathcal{O}=\{0,1\}$, and the size of their dwelling, denoted by $h>0$. Households who cannot own or rent or decide not to in order to increase their non-housing consumption, become homeless. We expand on the homelessness component in Section \ref{sec:cohabit_homeless}.

Households who purchase a house of size $h$ pay $p h$, with the restriction that homeowners must purchase a house no smaller than $\underline{h}_{o}$. Buying and selling houses involve proportional transaction costs, $\tau_{b} p h$ and $\tau_{s} p h$, respectively. All housing is subject to a property tax rate $\tau_{p}$ and depreciates at a rate $\delta$ per period. 

Renters must consume at least $0<\underline{h}_{r}<\underline{h}_{o}$ \citep{gervais,sommeretal-2013,sommersullivan-2017,kaplanetal-2019}, which we calibrate to match the share of renters spending more than 30\% of their pre-tax income on housing. This parameter captures the extent to which some renters prefer to consume less housing services than $\underline{h}_{r}$ but are prevented from doing so by exogenous constraints — such as natural limits, safety codes, or minimum size regulations imposed by municipalities.\footnote{See Section 9.5.4 of the Ontario Building Code: \url{https://www.ontario.ca/laws/regulation/120332}.} 

Controlled renters are also restricted to consuming at most $\overline{h}_c$ units of housing. Free-market renters, by contrast, can consume any housing unit, and their rents are denoted by $p_{r,1}$, which remain constant throughout their tenancy. All renters face a fixed cost of moving to a new rental unit, denoted by $\tau_r$, while working-age controlled renters additionally incur an explicit cost proportional to their labor income, $\varkappa\, y_{\ell,j}(\chi,e)$. This reduced-form cost is meant to capture time-use distortions under rent control — such as longer commuting times due to worse locations — without introducing an explicit heavy spatial structure. We explain how we calibrate $\varkappa$ in Section \ref{sec:internal_calibration}. 

\subsubsection{Tenancy-based rent control}

Under a tenancy-based rent control regime, renters can freely choose between the controlled and free rental markets, indicated by $f \in \mathcal{F}=\{0,1\}$, where $f=0$ denotes the controlled market and $f=1$ denotes the free market. The rent payment in the controlled market declines over time as follows:
\begin{equation}\label{eq:rent_in_controlled_market}
\frac{\tilde{p}_{r,0}}{(1+\kappa)^{\iota}}.
\end{equation} 

\noindent $\tilde{p}_{r,0}$ is the initial rent paid by a new tenant, $\kappa$ is the rate of rent decline for each additional period the tenant remains in the same unit, and $\iota$ is the number of periods the tenant has lived in the unit. 

Figure \ref{fig:diagram_rent_path} qualitatively illustrates free-market rents and controlled rents under tenancy control. The downward-sloping blue line shows that rents in the controlled market begin at $\tilde{p}_{r,0}$ and decline with tenancy duration, as per Equation \eqref{eq:rent_in_controlled_market}. The horizontal red line represents the free-market rent $p_{r,1}$, which is constant irrespective of tenancy duration, assuming the renter has not opted to allow their unit to deteriorate. 

When both markets coexist, the free-market rent must exceed the average controlled rent; otherwise, there would be no incentive to supply free-market rentals. In contrast, in equilibrium, the free-market rent $p_{r,1}$ must be lower than the initial rent in the controlled market $\tilde{p}_{r,0}$. What explains this equilibrium outcome?

If the initial rent in the controlled market were lower than the free-market rent $p_{r,1}$, too many renters would choose the controlled market, outstripping supply. To achieve equilibrium, therefore, renters must compete for controlled units, which results in an "entry premium" that clears the controlled market via self-selection. For instance, renters who expect to remain a short time in the unit would choose to rent in the free market to avoid this "entry premium". In contrast, renters who expect to stay longer in the unit would be more willing to pay the "entry premium" in exchange for longer-term discounts. Because landlords cannot observe a renter's expected duration when they move in, they must set a uniform entry rent, and the market clears via self-selection.\footnote{While tenancy control could, in principle, induce landlords to discriminate among prospective tenants when setting entry rents, we find no significant impact of observable characteristics — such as age or income — on rents paid by new tenants (see Section \ref{section:appendix_discrimination} in the Appendix). This finding supports the view that large landlords — who dominate the supply of controlled units — rely on standardized pricing, lack information about expected tenancy duration of specific prospective tenants, and face legal constraints that limit this type of price discrimination. See \citet{be-2000} for further discussion.}

The "entry premium" paid in the controlled market ensures that shortages do not arise. While this point has been highlighted previously by \citet{arnott-1995}, our framework formalizes it and shows that it also means that new movers in the controlled sector must pay rents higher than free-market rents, a point \citet{arnott-1995} does not make. Because the "entry premium" allocates controlled units to those with the highest willingness to pay at the time they move into the unit, the initial (ex-ante) allocation of controlled units is relatively efficient. However, because circumstances can change over time and controlled renters might "overstay their welcome" due to the lock-in effects in our model, allocation can still be inefficient ex-post. 

Using 2016 Census data for Toronto, we estimate that new tenants in the controlled market pay a 5.9\% "entry premium" over free-market rents.\footnote{See Section \ref{section:appendix_measuring_premium} in the Appendix.} This result aligns with prior empirical studies \citep{nagy-1997,early-2000,cjq-2023}. 

\subsubsection{Unit-based rent control}

Rents under a unit-based rent control system, denoted as $\hat{p}_{r,0}$, are regulated at the unit level below the free-market rate $p_{r,1}$, remain constant throughout a tenancy, and cannot be increased by landlords upon tenant turnover. Thus, new tenants receive immediate discounts relative to free-market rents. 

Because there is no price mechanism to allocate controlled units under unit control, shortages arise, and thus a random draw $\eta$ determines which households may gain access to controlled units. This draw is i.i.d. across time and all households, whether they are homeowners, renters, or homeless. With probability $\pi$, a household draws a favorable shock ($\eta=1$) and can rent a controlled unit, staying as long as they wish. However, they can turn this option down, and opt for buying a home, renting in the free market, or becoming homeless. Upon leaving their unit, they must draw a new favorable shock to be able to re-enter the controlled market. With probability $1-\pi$, the draw is unfavorable ($\eta=0$), and the household is excluded from the controlled market for that period. Those households can still rent in the free market, buy a home, or face homelessness. 

This highlights two key tradeoffs between unit control and tenancy control. The first concerns allocative efficiency. Because prices allocate controlled units under tenancy control but lotteries allocate them under unit control, the allocation at the time of entry is likely to be more efficient under tenancy control than under unit control. However, the allocation under tenancy control may become less efficient over time: tenure-dependent rent discounts strengthen lock-in effects over the course of a tenancy relative to unit control, potentially keeping renters in units long after their life circumstances have changed.

The second tradeoff concerns the timing of benefits. Under tenancy control, tenants must forgo current consumption in exchange for the prospect of future affordability gains as rent discounts accumulate with tenure. In contrast, unit control provides an immediate rent discount relative to free-market rents for households that win the lottery and secure a controlled unit, but may exclude households with very high willingness to pay. We discuss the third tradeoff right below, in Section \ref{sec:eviction}.

Our simulations in Section \ref{sec:non_targeted_statistics} show that unit control results in fewer controlled units being allocated to low-income households, but they also show that the allocation of controlled rentals under unit control is far from random: While the opportunity to rent a controlled unit — determined by lotteries — is random, the decision to actually rent a controlled unit and remain in it longer is not. Since low-income households are more likely to opt for a low-quality controlled unit and stay longer in the same unit, they ultimately still receive a disproportionately larger share of controlled units in equilibrium.

\subsubsection{Rent control and evictions}\label{sec:eviction}

A common argument against tenancy control is that it motivates landlords to evict long-term tenants in order to secure higher rents from new tenants. This incentive, however, does not arise in free-market or unit-control settings, where landlords cannot squeeze higher rents from new tenants relative to what incumbents already pay. To capture this mechanism, our model specifies an eviction probability that depends solely on the rent gap between what the landlord could charge a new tenant and what the current tenant pays.

Formally, renters face no eviction risk if they reside in free-market or unit-controlled rentals. Under tenancy control, however, the probability of eviction rises with the rent increase a landlord can gain by replacing the incumbent tenant with a new tenant. We specify the eviction probability $\epsilon$ as:
\begin{equation}\label{eq:eviction}
\epsilon(\iota)=
\begin{cases}
0 & \text{if not a renter or renting a free-market or unit-controlled unit} \\ \epsilon_1 \underbrace{\left(1-\frac{1}{\left(1 + \kappa \right)^{\iota}}\right)}_{\text{Depth of discount}} & \text{if renting a tenancy-controlled unit}.
\end{cases}
\end{equation}

\noindent Here, $\epsilon_1$ is a scale parameter that determines how strongly the rent gap influences eviction risk, $\epsilon(\iota)$. We calibrate $\epsilon_1$ to match the observed average eviction rate for reasons other than non-payment of rent, which we interpret as displacement driven by landlords' desire to obtain higher rents from new tenants.

Renters who are evicted from the tenancy-controlled market re-enter the next period as free-market renters, with their housing market state reset to $f=1$.

Because eviction associated with this mechanism in our Toronto data appears to be relatively uncommon — likely due to strong tenant-protection laws — this channel has limited quantitative bite in our experiments. Nevertheless, since it could be more consequential in institutional environments with weaker tenant protections, we decided to include it in the model.

\subsubsection{Homelessness}\label{sec:cohabit_homeless}

Because rental housing has a binding minimum size, some households may be unable to afford the conventional rental market or may choose not to participate if doing so would leave their non-housing consumption too low. Those households face homelessness $\left(b \in \{0,1\}\right)$. As in \citet{az-2024}, homeless households receive reduced housing services $\underline{h}_{h}$ free of charge, where $0<\underline{h}_{h}<\underline{h}_{r}$. However, as in \citet{abramson-2024}, the city incurs a cost $\theta_{b}$ per homeless household. 

For convenience, for households who are homeless, we set $o=0$ and $f=1$ when we define the dynamic program later.

\subsection{Mortgages}\label{subsection:mortgages}

In addition to rent control, the structure of mortgage contracts also influences households' tenure decisions, which in turn affect rental demand, equilibrium rents, and welfare results. Modeling mortgages as realistic long-term contracts is therefore important. Denoted as $m$, mortgages are long-term contracts with origination costs, prepayment penalties,\footnote{Prepayment penalties in the U.S. are common during the first 1-3 years of a mortgage and less so in later years.} and loan-to-value (LTV) and payment-to-income (PTI) constraints at origination only. We set the mortgage interest rate at $r_{m}=r+\Delta_m$, where $r$ is the return on financial assets, and $\Delta_m$ is the mortgage debt interest rate premium \citep{sommeretal-2013,sommersullivan-2017,kaplanetal-2019}. Consistent with Canadian tax policy, mortgage interest payments cannot be deducted from taxable income.

To obtain a new mortgage, households must pay a fixed origination cost $\omega$. A mortgage can only be originated if the LTV ratio does not exceed a threshold $\lambda_{1}(ph)<1$, which is a function of the house value $ph$, at the time of purchase:
\begin{equation}\label{eq:LTV}
m \leq \lambda_{1}(ph) \times \,p h
\end{equation}

\noindent The PTI constraint caps mortgage debt based on the borrower's pre-tax income at the time of origination: 
\begin{equation}\label{eq:PTI}
\frac{\tau_{p}ph+(r_{m}+\upsilon)m}{y_{\ell,j}(\chi,e)+y_{k}(a)} \leq \lambda_{2}.
\end{equation}

\noindent $\upsilon$ represents the average annual fraction repaid over the mortgage term, and $\lambda_{2}$ is the maximum PTI ratio permitted. This constraint ensures that the sum of property taxes, interest payments, and average principal repayments does not exceed a fixed share $\lambda_{2}$ of the borrower's pre-tax labor and capital income at the time of origination.

Each period, mortgage holders must allocate at least a fraction $\upsilon$ of the outstanding balance for repayment. We consider borrowers paying less than $\upsilon \times m$, as ones who must originate a new mortgage. Therefore, they must incur the fixed cost $\omega$ and adhere to constraints \eqref{eq:LTV}--\eqref{eq:PTI} once again.

\subsection{Summary of heterogeneity}

Households are heterogeneous along several dimensions:

\begin{itemize}
    \item Financial assets, $a\in \mathcal{A}=\mathbb{R}_{+}$.
    \item Housing tenure status, $o \in \mathcal{O}=\{0,1\}$.
    \item Opportunity to rent in the controlled market, in the case of unit control, $\eta \in \mathcal{U}=\{0,1\}$.
    \item Rental market type for renters, $f \in \mathcal{F}$.
    \item Number of periods a renter has lived in their unit, $\iota \in \mathcal{I}=\mathbb{N}$.
    \item Housing size, $h\in \mathcal{H}=\mathbb{R}_{++}$.
    \item Mortgage debt, $m \in \mathcal{M}=\mathbb{R}_{+}$.
    \item Idiosyncratic labor fixed effect $\chi \in \mathcal{X}=\mathbb{R}_{++}$.
    \item Stochastic labor productivity shock, $e \in \mathcal{E} = \mathbb{R}_{++}$.
\end{itemize}
 
\noindent We denote the measure of age-$j$ households over endowments by $\Psi_{j}(a,o,\eta,f,\iota,h,m,\chi,e)$. For brevity, we represent the state of a household whose age is $j$ using $s$.

\subsection{Dynamic program}\label{subsection:household_optimization}
Given its state variables, a household must make several decisions: saving ($a'$), housing tenure ($o'$), whether to be homeless ($b'$), and housing consumption ($h'$). 

Under tenancy control, renters also decide whether to rent in the controlled or free market ($f'$). However, we use $f'_z$ to denote the household’s realized rental-market state in the next period, distinguishing it from the chosen status $f'$. This is necessary because, in the event of an eviction, controlled renters are forced into $f'_z = 1$ regardless of their choice $f'$. Renters in the controlled market also choose whether to extend their tenancy by one more period ($\iota'$). We again distinguish between $\iota'$ and the realized $\iota'_z$ in the event of an eviction, in which case $\iota'_z=1$.

Homeowners further choose whether to originate a new mortgage ($n$) and their mortgage debt ($m'$). At the terminal age $J$, households also decide how much bequest to leave behind ($q'$).

To economize on space and avoid presenting multiple dynamic problems, we embed the owner, renter, and homeless dynamic problems within a single problem and represent tenancy control and unit control through a unified budget constraint. We define the rent control regime by $\alpha \in \{0,1\}$, where $\alpha=0$ corresponds to tenancy control and $\alpha=1$ to unit control.

The household's dynamic optimization problem is given by:
\begin{equation}\label{eq:vf}
\begin{aligned}
V_{j}(s) 
= \max_{\substack{\{c,\, a',\, o',\, b',\, f', \\ \iota',\, h',\, n,\, m',\, q'\}}} \Big\{
& \; u_{j}\big( c,\, h' \big)
+ \mathbbm{1}_{\{ j = J \}}\, \omega_{B}(q')+ \\
& \mathbbm{1}_{\{ j < J \}}\, \beta 
\int_{\mathcal{E}} \Big[
(1 - \epsilon(\iota'))\, V_{j+1}\big( s'(f'_z=f',\, \iota'_z=\iota') \big)d\tilde{F}(e'\mid e, b') +\\ & \quad \epsilon(\iota')\, V_{j+1}\big( s'( f'_z = 1,\, \iota'_z = 1 ) \big)d\tilde{F}(e' \mid e, b') \,  \Big]
\Big\},
\end{aligned}
\end{equation}

\noindent subject to the following constraints:
\begin{flalign}
  & c(1+\tau_{c})+a'+o r_{m} m+(1-b')(1-o')\left[(1-f')\left[ \frac{\left((1-\alpha)\tilde{p}_{r,0}+\alpha \hat{p}_{r,0}\right)} {(1+\kappa)^{\iota'-1}}h'+\mathbbm{1}_{\{j< J_{R}\}}\varkappa y_{\ell,j}(\chi,e)\right]+f'p_{r,1}h'\right]+  \label{eq:bc} \\ & \quad \mathbbm{1}_{\left\{[(o=1 \land o'=0)] \vee \left\{(o=0 \land o'=0) \land [(f' \neq f) \vee (h'\neq h)]\right\}\right\}}\tau_r+o\left(\tau_{p}+\delta+\mathbbm{1}_{\left\{o'=0 \vee h'\neq h\right\}}\tau_{s}\right)ph + o'(1+\tau_{b}\mathbbm{1}_{\left\{o=0 \vee h'\neq h\right\}})ph' + \notag \\ & \quad \omega n = y_{\ell,j}(\chi,e)+y_{k}(a)-y^{\tau}_{j}(\chi,e,a)+ a+o(ph-m)+ o'm'+\vartheta(\chi,e)+T(s)+T_{fed}(s) \notag \\
    & a' \geq 0 \label{eq:financial_wealth} \\
    & h' \geq o'\underline{h}_{o}+(1-o')\underline{h}_{r} \quad \text{if $b'=0$} \label{eq:min_house_size} \\
    & h' \leq \overline{h}_{c} \quad \text{if $f'=0$} \label{eq:max_house_size} \\
    & h'=\underline{h}_{h} \quad \text{if $b'=1$} \label{eq:min_house_size_homeless} \\
    & f' = 1 \quad \text{if $\eta=0$ and $\alpha=1$}\label{eq:free_market} \\
    & n \in \left\{0,o'\mathbbm{1}_{\left\{j<J_{R}\right\}}\right\} \label{eq:origination} \\
    & nm' \leq \lambda_{1}(ph') p h' \label{eq:LTV_1} \\
    & n (\tau_{p}ph'+(r_{m}+\upsilon)m')\leq \lambda_{2}(y_{\ell,j}(\chi,e)+y_{k}(a))\label{eq:PTI_1} \\
    & (1-n)m'\leq (1-\upsilon)m \label{eq:settle_debt} \\
    & q' = a' + o' \left[p h'(1-\tau_p-\delta-\tau_s)-(1+r_m)m'\right] \label{eq:networth}
\end{flalign} 

\noindent The household's utility function comprises three components: a flow utility derived from non-housing and housing consumption, a warm-glow component stemming from bequests left behind by deceased households, and a discounted value of future utility which depends on whether the household was evicted or not.\footnote{The dynamic program is written to allow for labor productivity shocks across all ages, facilitating a unified representation of decision-making for both retirees and working-age individuals.}

Equation \eqref{eq:bc} defines the household's budget constraint. The left-hand side includes expenditures (in order of appearance): consumption and consumption taxes, savings, mortgage interest payments, rental costs in controlled and free markets, time-use costs for working-age controlled renters, moving costs renters entering a new unit, property taxes, housing depreciation, housing transaction costs (buying/selling), and mortgage origination costs. The right-hand side comprises resources available for expenditure (in order of appearance): after-tax labor and capital incomes, financial assets, housing equity, mortgage funds raised, bequests contingent upon the household's labor productivity, lump-sum transfers from the rental and construction companies that depend on the household's state, and lump-sum transfers from the federal government that may also be conditional on the state of the household. In our simulations, we rebate profits exclusively to households in the top of the wealth distribution, consistent with empirical evidence on the concentration of landlordship at the top of the wealth distribution.

The remaining constraints specify behavioral and institutional limits: 
\begin{itemize}
    \item Constraint \eqref{eq:financial_wealth}: Households cannot borrow without collateral (no unsecured debt).
    \item Constraint \eqref{eq:min_house_size}: All owners must consume at least $\underline{h}_{o}$ units of housing, and all renters must consume at least $\underline{h}_{r}$.
    \item Constraint \eqref{eq:max_house_size}: Controlled renters' housing consumption cannot exceed $\overline{h}_{c}$.
    \item Constraint \eqref{eq:min_house_size_homeless}: Households experiencing homelessness must consume exactly $\underline{h}_{h}$ unit of housing services.
    \item Constraint \eqref{eq:free_market}: Under unit control, households cannot rent in the controlled market if they draw a bad shock ($\eta=0$). 
    \item Constraint \eqref{eq:origination}: Only working-age households may originate new mortgages.
    \item Constraints \eqref{eq:LTV_1} and \eqref{eq:PTI_1}: LTV and PTI constraints at origination.
    \item Constraint \eqref{eq:settle_debt}: Mortgage holders must repay at least a minimum fraction $\upsilon$ of their principal each period.
    \item Constraint \eqref{eq:networth}: Bequests equal net financial and net housing wealth at death, accounting for the mortgage, property taxes, depreciation, and selling costs.
\end{itemize} 

Finally, we denote policy functions for households as follows: consumption $c_{j}(s)$, savings $a'_{j}(s)$, homelessness $b'_{j}(s)$, housing tenure $o'_{j}(s)$, rental market choice $f'_{j}(s)$, the number of periods lived in the rental unit $\iota'_{j}(s)$, housing size $h'_{j}(s)$, mortgage debt $m'_{j}(s)$, mortgage origination $n_{j}(s)$, and bequests $q'_{j}(s)$.

\subsection{Rental company with partial exemptions}\label{sec:rental_company}

\subsubsection{Overview}

A central concern with any rent control policy is its impact on the supply of rental housing. To address this, we develop a rental supply framework that extends and formalizes the partial-exemption logic of \citet{fallis-1985} within a general equilibrium setting.

Before introducing the formal optimization problem and its solution, we describe the rental company’s decision problem in plain terms. The rental company chooses how many housing units to operate as rentals, taking housing prices, rents, and regulations as given. A key institutional feature of rent control is that only a subset of the housing stock — typically newer buildings — is legally exempt from regulation. Even within this exempt stock, however, only a fraction can be converted into rental use.\footnote{In practice, legal restrictions such as condominium bylaws and cooperative rules, as well as physical limitations, cap the fraction of exempt housing that can be used as rental units. We capture these constraints by imposing an exogenous upper bound on rental utilization.} The remaining exempt units are absorbed by owner-occupiers rather than left vacant.

This utilization constraint places an upper bound on the quantity of rental housing that can be supplied at free-market rents using exempt stock alone. Importantly, this feature makes partial exemptions economically meaningful: without it, even a small exempt segment would suffice to replicate the unconstrained rental allocation, rendering rent control irrelevant in equilibrium.

As long as free-market rents exceed marginal costs, the rental company supplies rentals from the exempt stock up to this utilization limit. If the exempt stock is sufficiently large, this capacity alone may be enough to reach the unconstrained competitive rental quantity. When exemptions are limited, however, the firm may find it profitable to expand rental supply into the controlled segment. It does so only if the expected controlled rent covers the marginal cost of operating an additional controlled rental unit. As a result, controlled units are never loss-making and are not cross-subsidized by free-market rentals: they appear in equilibrium only when the utilization constraint on exempt housing binds and supplying regulated units remains privately profitable.

A critical ingredient of the rental supply block is the presence of convex management costs. These costs generate a realistically upward-sloping rental supply curve, consistent with empirical evidence that long-run rental supply is far from perfectly elastic. Absent convex costs, rental supply would be perfectly elastic at the user cost of housing, implying that any binding rent control would mechanically eliminate rental supply and could never be welfare-improving. Convex costs therefore play a central role in allowing rent control to affect supply without collapsing it entirely, making a quantitative welfare analysis meaningful.

This modeling choice contrasts with frameworks in which households endogenously choose whether to become landlords, which typically generate very large rental supply elasticities \citep{rs-2023}. Even recent quantitative models that incorporate realistic landlord frictions — such as lumpy adjustment and borrowing constraints — find steady-state elasticities around 3.5 \citep{araujo2024}, still well above empirical estimates for large urban markets.

\subsubsection{Formalizing the model}

We now formalize the rental company’s problem. We first set up the optimization problem and then explain how it is solved.

\subsubsection*{Setting up the optimization problem}

The rental company is a price taker and supplies rental housing in two regulatory segments:

\begin{itemize}
    \item \textbf{Free-market segment} ($i=1$): A subset of the housing stock, denoted $H'_1$, is exempt from rent control and earns the free-market rent $p_{r,1}$. The size of the exempt stock is determined by regulation and taken as given. However, only a fraction $\phi \in (0,1)$ of this stock can be operated as rentals. The remaining $(1-\phi)H'_1$ units are occupied by owner-occupiers. We calibrate $\phi$ by estimating rental utilization rates in Calgary and Edmonton — two large Canadian cities without rent control — and weighting these rates by Toronto’s housing composition, yielding an estimated utilization rate of $46.8\%$ for exempt housing. In other words, at most $46.8\%$ of any exempt housing stock can be converted into rentals.
    
    \item \textbf{Controlled segment} ($i=0$): The remainder of the housing stock, $H'_0 = H' - H'_1$, is subject to rent control when rented and earns the expected controlled rent $\overline{p}_{r,0}$.
\end{itemize}

To supply a total of $X' = X'_0 + X'_1$ rental units, the firm incurs convex management costs
\[C(X') = \frac{\theta_1}{\theta_2} X'^{\theta_2}, \qquad \theta_2 > 1,\]
as in \citet{garrigaeetal-2009}. Each rental unit depreciates at rate $\delta$ and is subject to a property tax $\tau_p$. Housing can be purchased at price $p$.

Conditional on rents, prices, and the size of the exempt stock, the rental company solves
\begin{align}\label{eq:rental_optimization}
V(X) &= \max_{\{X'_0, X'_1\}} \Big\{\overline{p}_{r,0} X'_0 + p_{r,1} X'_1- \frac{\theta_1}{\theta_2} (X')^{\theta_2}- p X' + p(1-\delta-\tau_p) X+ \frac{1}{1+r} V(X')\Big\},
\end{align}
subject to regulatory constraints and
\[X' = X'_0 + X'_1.\]

\noindent The implied convex marginal cost of rental supply is
\begin{align}\label{eq:rental_MC}
MC(X') = \theta_1 X'^{\theta_2 - 1}+ p\left(\frac{r + \delta + \tau_p}{1+r}\right).
\end{align}

Period profits are given by
\[\Pi_r=\underbrace{\overline{p}_{r,0} X'_0 + p_{r,1} X'_1}_{\text{revenues}}-\underbrace{\left[\frac{\theta_1}{\theta_2} (X')^{\theta_2}+ p(\delta + \tau_p) X'\right]}_{\text{costs}}.\]

\subsubsection*{Solving the optimization problem}

To characterize optimal rental supply, we solve the firm’s problem by backward induction. This solution method reflects the fact that, with convex management costs, the profitability of supplying regulated units depends on whether the utilization constraint on exempt housing binds in equilibrium. Although free-market rentals appear first in the description of equilibrium outcomes, the firm’s optimization problem is naturally solved in reverse order, because the decision to supply controlled units determines the marginal cost schedule relevant for subsequent free-market supply.

Let $X'^*$ denote the unconstrained competitive rental quantity that would be supplied if a sufficiently large share of the housing stock were exempt from regulation. This benchmark satisfies
\[MC(X'^*) = p_{r,1}^*,\]
\noindent where $p_{r,1}^*$ is the competitive free-market rent absent regulation.

When the exempt stock is limited and the utilization constraint binds, the firm optimally supplies free-market rentals up to the utilization limit,
\begin{equation}\label{eq:uncontrolled_optimization}
X'_1 = \phi H'_1 < X'^*,
\end{equation}
since free-market rentals yield strictly positive surplus at the margin up to this cap.

Anticipating that free-market rentals will be supplied up to this limit, the firm then determines whether it is profitable to expand rental supply in the controlled segment. Importantly, this decision is made \emph{before} free-market rentals are supplied. As a result, the relevant marginal cost for controlled supply is the marginal cost of increasing rental activity starting from zero rentals and evaluated at $X'_0$ — the supply of controlled rentals — rather than the marginal cost starting from the supply of free-market rentals $X'_1$, and evaluated at the total quantity $X'=X'_0 + X'_1$.

Controlled supply is therefore chosen such that the expected controlled rent equals the marginal cost of supplying an additional controlled unit:
\begin{equation}\label{eq:controlled_optimization}
X'_0=\min\left\{\max\left\{X'^*_0, 0\right\},\; X'^* - X'_1
\right\},\qquad \text{where} \qquad MC(X'^*_0) = \overline{p}_{r,0}.
\end{equation}

\noindent It is important to draw the reader's attention to the fact that the optimal controlled rental supply given by Equation \eqref{eq:controlled_optimization} is bounded above by the gap between the unconstrained efficient rental quantity and the free-market rentals supplied from the exempt stock. If the controlled supply implied by this equation exceeds this gap, the firm optimally reduces controlled supply to exactly fill the remaining gap, since supplying additional controlled units would push total rental supply beyond the efficient level and raise marginal costs above both controlled and free-market rents. 

In contrast, if expected controlled rents are sufficiently low, the firm supplies fewer controlled units than would be required to reach the unconstrained competitive quantity. In this case, total rental supply falls short of $X'^*$, and rent control restricts rental housing supply in equilibrium.

\subsection{Housing construction company}\label{sec:construction_company}

A competitive construction firm earns revenue $pI$ from producing $I$ units of new housing. Its total cost of production is given by:

\begin{equation*}\label{eq:total_cost_structure}
TC(I)=\frac{\varsigma_{1}}{\varsigma_{2}}I^{\varsigma_{2}},
\end{equation*}

\noindent where $\varsigma_{1}$ dictates the average cost of housing investment, and $\varsigma_{2}$ controls the responsiveness of investment to house prices. The firm chooses $I$ to maximize profits:
\begin{equation*}
\label{eq:construction_company_profits}
\Pi_h=\max_{I}\left\{pI-TC(I)\right\}.
\end{equation*}

\noindent The first-order condition for this problem implies the following standard positive relationship between housing investment and the price of housing:
\begin{equation}
\label{eq:construction_company_first_order_condition}
I^{\star}=\left(\frac{p}{\varsigma_{1}}\right)^{\frac{1}{\varsigma_{2}-1}}.
\end{equation}

\subsection{Aggregation}\label{subsection:aggregation}

Rental and construction company profits must equal the total lump-sum transfers from rental and construction companies to households, leading to the following constraint:
\begin{align}\label{eq:rebates}
   (\Pi_h+\Pi_r)=\sum_{j=1}^J \int_\mathcal{S} T(s)\mathrm{d}\Psi_j(s).
\end{align}

As outlined in Section \ref{sec:brief_overview}, our model incorporates both local and federal governments. The local government collects property taxes, and provides local public goods to residents, which do not enter into households' utility, as well as shelters to the homeless. Therefore, the budget constraint of the local government is as follows:
\begin{align}\label{eq:gbc_1}
  G_{local} +\sum_{j=1}^J \int_\mathcal{S} \theta_{b} b'_{j}(s)\mathrm{d}\Psi_j(s) = \tau_{p} p H.
\end{align}

\noindent In calibrating the model to match the existing tax code in Toronto, we set the level of consumption of local public goods, denoted as $G_{local}$, to ensure budget balance. In all policy experiments, we hold $G_{local}$ constant and adjust the property tax rate $\tau_{p}$ to satisfy the budget constraint.

The federal government collects income and consumption taxes and provides federal public goods $G_{fed}$, which also do not enter households' utility, as well as pensions to retirees and lump-sum transfers to households $T_{fed}(s)$. As in the case of the local government, we set the level of $G_{fed}$ to ensure budget balance in our benchmark, and in our experiments, we hold $G_{fed}$ at its benchmark level, and to balance the budget, we adjust either $\tau_\ell$, which determines the average labor income tax, or the lump-sum transfers, $\{T_{fed}(s)\}_{s\in \mathcal{S}}$. Thus, we require that:
\begin{align}\label{eq:gbc_2}
   G_{fed}+\sum_{j=J_R}^J \int_\mathcal{S} y_{R}(s)\mathrm{d}\Psi_j(s)+\sum_{j=1}^J \int_\mathcal{S} T_{fed}(s) \mathrm{d}\Psi_j(s)=\sum_{j=1}^J \int_\mathcal{S} \left(y^{\tau}_{j}(s)+\tau_{c}c_{j}(s)\right)\mathrm{d}\Psi_j(s) 
\end{align}

Before presenting the remaining market-clearing conditions, we consider measurable sets $S \subseteq \mathcal{S}$. For notational convenience, we sometimes work with product sets of the form
\[
S = A \times O \times U \times F \times I \times H \times M \times X \times E,
\]

\noindent where $A \subseteq \mathcal{A}$, $O \subseteq \mathcal{O}$, $U \subseteq \mathcal{U}$, $F \subseteq \mathcal{F}$, $I \subseteq \mathcal{I}$, $H \subseteq \mathcal{H}$, $M \subseteq \mathcal{M}$, $X \subseteq \mathcal{X}$, and $E \subseteq \mathcal{E}$, respectively.

We require that aggregate housing demand equals aggregate housing supply:
\begin{equation}
\label{eq:housing_clearing_condition}
  \sum_{j=1}^J \int_\mathcal{S} (1-b'_{j}(s))h'_{j}(s) \mathrm{d}\Psi_j(s) = H',
\end{equation}

\noindent where the law of motion of the housing stock is: 
\begin{equation}\label{eq:housing_stock_lom}
H'=I+(1-\delta)H.
\end{equation}

\noindent Equation \eqref{eq:housing_stock_lom} is a standard law of motion for housing, stating that the next period's housing stock equals new investment plus the depreciated stock of housing. 

We also impose market-clearing conditions for rental housing. In each rental market, aggregate rental demand equals aggregate rental supply:
\begin{equation}\label{eq:rental_housing_clearing_condition}
  \sum_{j=1}^J \int_\mathcal{S} (1-o'_{j}(s))\mathbbm{1}_{\left\{f'_{j}(s)=i\right\}}(1-b'_{j}(s))h'_{j}(s) \mathrm{d}\Psi_j(s) = X'_{i} \quad i \in \{0,1\}.
\end{equation}

The distributions of households evolve according to the following laws of motion:
\begin{equation}\label{eq:lom1}
\Psi_{1}(S) = \int_{\mathcal{X} \times \mathcal{E}} \mathbbm{1}_{\big\{ (0,0,1,0,0,0,0,X,E) \in S \big\}} d\hat{F}(X)\, d\tilde{F}(E), \quad \text{under tenancy control.}
\end{equation}

\begin{equation}\label{eq:lom1_2}
\begin{aligned}
\Psi_{1}(S)={}&\;\pi \int_{\mathcal{X}\times\mathcal{E}}
\mathbbm{1}_{\{(0,0,1,1,1,1,0,X,E)\in S\}}\, d\hat F(X)\, d\tilde F(E)\\&\;+(1-\pi)\int_{\mathcal{X}\times\mathcal{E}}
\mathbbm{1}_{\{(0,0,0,1,1,1,0,X,E)\in S\}}\, d\hat F(X)\, d\tilde F(E),\qquad \text{under unit control.}
\end{aligned}
\end{equation}

\noindent Under unit control, newborn households draw the opportunity shock $\eta$ at entry from the same distribution as incumbent households.

\begin{equation}\label{eq:lom2}
\begin{aligned}
\Psi_{j+1}(S)={}&\int_{\mathcal{S}}\Bigg[\int_{\mathcal{E}}
\Big([1-\epsilon(\iota'_j(s))]\,\Xi^{NE}_j(s,E';S)
+\epsilon(\iota'_j(s))\,\Xi^{E}_j(s,E';S)\Big)
\\[-0.2em]
&\hspace{4em}\times d\tilde F\!\left(E' \mid E,b'_j(s)\right)\Bigg]\,L_j(s)\, d\Psi_j(s),\qquad j=1,\ldots,J-1.
\end{aligned}
\end{equation}

\noindent Let $\eta\in\mathcal{U}=\{0,1\}$ denote the opportunity shock that governs access to a controlled unit under unit control.  Under unit control,
\[
\Pr(\eta=1)=\pi,
\]
where $\pi$ adjusts to clear the controlled market. Under tenancy control, $\eta$ is irrelevant, so for convenience we set it to $1$. We include it only to keep a unified state-space notation across regimes. Eviction is governed by $\epsilon(\iota'_j(s))$ and corresponds to the transition $s'^{\,E}_j$, defined shortly below; the shock $\eta$ is a separate opportunity shock relevant only under unit control.

\noindent Define the implied indicator terms (i.e., the expectation over $\eta'$) as
\begin{equation}\label{eq:Xi_def}
\begin{aligned}
\Xi^{NE}_j(s,E';S)&=\begin{cases}
\mathbbm{1}\!\left\{ s'^{\,NE}_j(s,1,E') \in S \right\}
& \text{under tenancy control},\\[0.35em]
(1-\pi)\,\mathbbm{1}\!\left\{ s'^{\,NE}_j(s,0,E') \in S \right\}
+\pi\,\mathbbm{1}\!\left\{ s'^{\,NE}_j(s,1,E') \in S \right\}
& \text{under unit control},
\end{cases}
\\[0.6em]
\Xi^{E}_j(s,E';S)&=\begin{cases}
\mathbbm{1}\!\left\{ s'^{\,E}_j(s,1,E') \in S \right\}
& \text{under tenancy control},\\[0.35em]
(1-\pi)\,\mathbbm{1}\!\left\{ s'^{\,E}_j(s,0,E') \in S \right\}
+\pi\,\mathbbm{1}\!\left\{ s'^{\,E}_j(s,1,E') \in S \right\}
& \text{under unit control}.
\end{cases}
\end{aligned}
\end{equation}

\noindent The next-period state under no eviction is
\begin{equation}\label{eq:sprime_no_eviction}
s'^{\,NE}_j(s,\eta',E')=\bigl(a'_j(s),\,o'_j(s),\,\eta',\,f'_j(s),\,\iota'_j(s),\,h'_j(s),\,m'_j(s),\,\chi,\,E'\bigr).
\end{equation}

\noindent Under eviction, the household is forced into the free market and its tenancy status resets, so
\begin{equation}\label{eq:sprime_eviction}
s'^{\,E}_j(s,\eta',E')=\bigl(a'_j(s),\,o'_j(s),\,\eta',\,1,\,1,\,h'_j(s),\,m'_j(s),\,\chi,\,E'\bigr).
\end{equation}

\noindent Further, we define the feasibility indicator $L_j(s)$ as
\begin{equation}\label{eq:Lj}
L_j(s)=\mathbbm{1}_{\bigl\{a'_j(s)\in \mathcal{A},\,o'_j(s)\in \mathcal{O},\,f'_j(s)\in \mathcal{F},\,\iota'_j(s)\in \mathcal{I},\,
h'_j(s)\in \mathcal{H},\,m'_j(s)\in \mathcal{M}\bigr\}}.
\end{equation}

Finally, total bequests from deceased households must equal total wealth endowments received by living households:
\begin{equation}\label{market_clearing:agg_bequests}
    \int_{\mathcal{S}} q'_{J}(s) \mathrm{d}\Psi_{J}(s) = \sum_{j=1}^J \int_{\mathcal{S}} \vartheta (\chi,e)\ \mathrm{d}\Psi_j(s).
\end{equation}

\noindent Following \citet{kkk-2021}, we choose the parameter $\vartheta$ to ensure that this condition holds in all our experiments. Consequently, variations in $\vartheta$ across experiments are contingent upon proportional changes in aggregate bequests.

\subsection{Equilibrium}
A stationary equilibrium in the model consists of: 
\begin{itemize}
   \item A sequence of aggregate prices and quantities:\\   
    $\left\{H', I, \alpha, \{X'_{i}\}_{i=0}^{1},    \overline{p}_{r,0}, p_{r,1}, \Pi_r, \Pi_h, \tilde{p}_{r,0}, \hat{p}_{r,0}, \pi, p, \tau_{\ell}, \overline{\tau}_{\ell}, \{T(s), T_{fed}(s)\}_{s \in \mathcal{S}}, \vartheta \right\}$;
    \item A sequence of value and policy functions: \\
    $\biggl\{V_{j}(\cdot),c_{j}(\cdot),a'_{j}(\cdot),o'_{j}(\cdot),b'_{j}(\cdot),f'_{j}(\cdot),\iota'_{j}(\cdot),h'_{j}(\cdot),m'_{j}(\cdot),n_{j}(\cdot),q'_{j}(\cdot)\biggl\}_{j=1}^J$; and
    \item A sequence of measures: 
     $\left\{\Psi_{j}(\cdot)\right\}_{j=1}^J$,
\end{itemize}
\noindent such that:
\begin{enumerate}
\item The policy and value functions solve the household's dynamic problem \eqref{eq:vf} subject to constraints \eqref{eq:bc}--\eqref{eq:networth};
\item Rental supply in both markets $\{X'_{i}\}^{1}_{i=0}$, solves the rental company's problem as defined in \eqref{eq:rental_optimization};
\item Housing investment $I$, solves the construction company's first-order condition \eqref{eq:construction_company_first_order_condition};
\item The lump-sum transfers from the rental and construction companies $\{T(s)\}_{s \in \mathcal{S}}$, satisfy constraint \eqref{eq:rebates};
\item The property tax $\tau_p$, balances the local government's budget \eqref{eq:gbc_1};
\item Either $\tau_{\ell}$ or $\{T_{fed}(s)\}_{s\in \mathcal{S}}$ balance the federal government's budget \eqref{eq:gbc_2};
\item The housing price $p$, clears the housing market
\eqref{eq:housing_clearing_condition}, subject to the housing stock law of motion \eqref{eq:housing_stock_lom};
\item The free-market rent $p_{r,1}$, clears the free rental market (\eqref{eq:rental_housing_clearing_condition}, when $i=1$); and
\begin{itemize}
    \item Under tenancy control, the initial controlled rent $\tilde{p}_{r,0}$ equates the aggregate demand for controlled rentals to the aggregate supply of controlled units (\eqref{eq:rental_housing_clearing_condition}, with $i=0$); or
    \item Under unit control, the probability $\pi$ that the random shock $\eta=1$, which determines whether a household has an opportunity to rent a controlled unit, clears the controlled market (\eqref{eq:rental_housing_clearing_condition}, when $i=0$);
\end{itemize}
\item The laws of motion for the distribution of households \eqref{eq:lom1}--\eqref{eq:lom2}, evolve in accordance with households' policy functions; and
\item $\vartheta$ equates aggregate bequests left by decedents to aggregate bequests provided to all living households, \eqref{market_clearing:agg_bequests}.
\end{enumerate}

\section{Calibration}\label{sec:calibration}

Since Toronto's rental market operates with tenancy control regulations, we calibrate our model accordingly ($\alpha=0$). Later we explore how the calibrated model changes when unit control is imposed instead ($\alpha=1$).

The calibration of our model proceeds in two steps. First, we calibrate specific parameters using data from Toronto and existing estimates from the literature. Second, we calibrate the remaining parameters to match certain model-generated statistics with their corresponding moments in our Toronto data. Table \ref{tab:calibration_external} summarizes the externally assigned parameters and their sources. Table \ref{tab:calibration_internal} lists the internally calibrated parameters along with the corresponding empirical targets. In Section \ref{sec:non_targeted_statistics}, we assess the model's performance on non-targeted dimensions.

In Section \ref{appendix:data} of the Online Appendix, we provide details on how we constructed key statistics from various datasets. 

\subsection{Assigned parameter values}\label{sec:external_calibration}

\subsubsection{Demographics, preferences, and income}
A model period is one year. Households enter the model at age 26, live up to age 85, and retire at age 66. This implies a retirement age of $J_{R}=41$ and a terminal age of $J=60$. 

We set the coefficient of relative risk aversion $\sigma$ to 2, and assign the equivalent household size by age cohort, $\{\xi_{j}\}_{j=1}^{J}$, using data from the 2016 Census. We scale household sizes by their square root, following standard practice in the literature. We determine the discount factor $\beta$, the utility weight on housing services $\gamma$, and the strength of the bequest motive $\mu$, in the second stage of our calibration.

We adopt the labor productivity process proposed by \citet{gkkoc}. Since the available Toronto data do not allow us to estimate the same set of productivity moments, we use their estimates to calibrate this process. The standard deviation of the fixed effect is set to $0.128$, and the stochastic labor productivity shocks follow an AR(1) process in logarithmic terms,

\begin{equation*}
  \log e' = \rho_{e} \log e + \epsilon_{e},\quad \epsilon_{e}\overset{\mathrm{i.i.d.}} \sim N(0,\sigma^2_{e}).
\end{equation*}

\noindent As in \citet{gkkoc}, we set the persistence parameter $\rho_{e}$ to 0.9. The standard deviation of the stochastic component of labor productivity $\sigma_{e}$, is calibrated internally. Newborns draw their stochastic productivity from the unconditional probability distribution associated with the ergodic distribution, which is endogenously generated by our productivity process. Consistent with \citet{az-2024}, we modify the stochastic process for homeless households by increasing the probability of moving to or remaining in the lowest productivity state by 20\%.  

We use 2016 Toronto Census data to estimate the deterministic lifecycle productivity component as:

\begin{equation*}
\zeta_j = 0.923+0.09\times j-1.61\times 10^{-3}\times j^{2}, \mbox{ for } j\leq J_{R},
\end{equation*}

\noindent where the coefficients are obtained from a regression of average household income for each age cohort, ranging from 26 to 64, on age and age squared. 

We determine pension benefits ${y_{R}(\chi,e)}$, based on the replacement rate of gross individual earnings in Canada estimated by \citet{OECD-pension}. Surviving households' endowments are proportional to their labor income: $\vartheta(\chi,e):=\vartheta \chi e$. We select the parameter $\vartheta$, governing surviving households' wealth endowments, to fulfill Equation \eqref{market_clearing:agg_bequests}.

\subsubsection{Taxes and interest rates}

We set the taxable portion of capital income $\nu$ to $0.5$, consistent with the Canadian tax code. Utilizing household-level data from the 2016 Census for Toronto, we estimate the parameter $\psi=0.9275$ in Equation \eqref{eq:income_tax}, which determines the progressivity of income taxes. We set the consumption tax rate $\tau_{c}$ to 13\% to reflect the Harmonized Sales Tax (HST) in Ontario. The interest rate $r$ is set to $2.77\%$, following the estimate for Canada by \citet{bruneau-etal-2018}.

\subsubsection{Housing}

We set the per-household public cost of homelessness $\theta_{b}$ to $82.3\%$ of average household income. We determine this parameter by dividing the City of Toronto's 2021 total spending on homelessness supports by the number of homeless individuals.\footnote{See \url{https://www.homelesshub.ca/community-profile/toronto} and \url{https://www.toronto.ca/legdocs/mmis/2021/ex/bgrd/backgroundfile-163804.pdf}. In the vast majority of cases, individuals experiencing homelessness are single adults. As a result, we think that individual-level data on homelessness is a reasonable proxy for having household-level data on homelessness.}

We set the depreciation rate of housing $\delta$ to $0.6\%$ based on estimates from \citet{rosenthal-2014}. The property tax rate $\tau_p$ is set to $0.44\%$, based on data from the City of Toronto and the Altus Group. Following \citet{gm-2003}, we set buying costs $\tau_{b}$ to $2\%$ of the purchase price, and selling costs $\tau_{s}$ to $6\%$ of the sale price. We set renters’ moving costs $\tau_r$ to 2.3\% of the benchmark's average labor income, corresponding to \$2,500 in 2016. This parameter captures expenses related to searching for a new unit, packing, transportation, and hiring movers. For homeless individuals, it also reflects additional barriers, including lower acceptance probabilities due to perceived homelessness and the difficulty of affording upfront large payments such as first and last month’s rent.

To generate a price elasticity of new housing supply of 0.89 as estimated for Toronto by \citet{piaxao-2021}, we set the cost curvature parameter $\varsigma_{2}$ to $2.12$. The construction company's cost function level, denoted as $\varsigma_1$, is a free parameter that we derive from the market clearing condition \eqref{eq:housing_clearing_condition} in the benchmark, and then hold fixed at its benchmark value across our experiments.

We set the largest house size a controlled renter can pick $\overline{h}_{c}$ to equal the housing size associated with the second housing grid point of our housing grid. This size is 81\% of the average free-market rental in the benchmark, and generates a gap in house sizes between the free market and the controlled market of 24.1\%, which happens to be exactly consistent with the estimates of the same measure by \citet{zhang-gou}.\footnote{Our housing size grid is $\{0.30, 0.45, 0.58, 0.75, 0.94, 1.37, 1.65\}$, where $0.30$ corresponds to the minimum rental size $\underline{h}_{r}$, $0.45$ to the maximum housing size available in the controlled market $\overline{h}_{c}$, and $0.94$ to the minimum owner-occupied size $\underline{h}_{o}$. Both $\underline{h}_{r}$ and $\underline{h}_{o}$ are chosen in the second stage of the calibration. Housing services for the homeless, $\underline{h}_{h}$ are set to $0.11$ also in the second stage of our calibration.} In our benchmark about 20\% of renters occupy free-market housing units larger than $\overline{h}_{c}$, and in our counterfactual experiments in which we decontrol the rental market, about 50\% of renters choose sizes larger than that.

We pick $\theta_1$, which determines the free-market rent, to generate a free-market-rent-to-house-price ratio of 5.82\%, based on our analyses utilizing 2006 Rental Market Survey (RMS) data, 2006 Toronto Real Estate Board data, 2006 Census data, and estimates from \citet{zhang-gou}. We share our full calculation in the Appendix (\ref{appendix:data}).

We calibrate $\phi=0.468$ using 2016 Census data for Toronto, Calgary, and Edmonton. Specifically, we calculate the rental utilization rates for condominiums and single-family homes in Calgary and Edmonton, two major Canadian cities without rent control. The averages of type-specific rates are then combined using Toronto’s actual distribution of housing types to obtain the final estimate. 

Furthermore, we set the parameter $\theta_2$, which governs the convexity of the marginal cost in the rental housing sector and determines the rental supply elasticity, to 3.71. When $\theta_2=1$, the marginal management cost is constant, implying a perfectly elastic rental supply. As $\theta_2$ increases above one, the marginal cost rises more steeply with the scale of rental supply, reducing supply elasticity. The value we pick for $\theta_2$ generates a rental supply elasticity that is equal to 1 in our benchmark. This value is somewhat lower than the 1.4 estimate for the overall U.S. rental market reported by \citet{rs-2023}, but it is consistent with evidence that supply elasticities are substantially lower in dense urban areas \citep{cy-1994,susin-2002,fack-2006,sy-2020}, which are the very ones where government normally considers rent control. Thus, we see our choice of elasticity of 1 as reflecting an upper bound.

\citet{zhang-gou} estimate that quality-adjusted controlled rents are on average 13.7\% below free-market rents. In our model, the rental company’s problem is naturally solved backward from the unconstrained efficient allocation, where at the efficient quantity, marginal cost equals the free-market rent. With a rental supply elasticity of 1, marginal cost falls proportionally as total rental quantity declines. Since exempt (free-market) units account for about 30\% of Toronto’s rental stock according to the 2016 Census (and more than 30\% of quality-adjusted rental quantity), the remaining 70\% of rentals must be supplied from the controlled market. Evaluating marginal cost at this point — when total rental supply is 70\% of the unconstrained level — implies that marginal cost of a controlled unit is at least 30\% below the free-market rent. However, because the expected controlled rent is only 13.7\% lower than the free-market rent, supplying controlled units remains privately profitable up to 70\% of the efficient quantity. As a result, in our benchmark calibration rent control does not reduce aggregate rental supply, reflecting the generous exemptions.

More generally, for rent control to reduce aggregate rental supply in this environment, the long-run rental supply elasticity must be sufficiently large so that marginal cost rises slowly with rental quantity. Let $q \in [0,1]$ denote total rental quantity as a fraction of the unconstrained efficient rental quantity (so $q=1$ corresponds to the efficient allocation). In the unconstrained allocation, marginal cost equals the free-market rent, i.e., $MC(1)=p_{r,1}$. 

Under constant-elasticity rental supply, marginal cost scales with quantity as
\[
\frac{MC(q)}{MC(1)} = q^{1/\varepsilon},
\]
where $\varepsilon$ is the long-run rental supply elasticity. In our setting, the firm supplies controlled units up to $q=0.70$ before supplying the remaining exempt units. Thus, a reduction in aggregate rental supply occurs only if the expected controlled rent, $0.863\,p_{r,1}$, falls below marginal cost at the margin when $q=0.70$, i.e.,
\[
0.863\,p_{r,1} < MC(0.70) = MC(1)\cdot 0.70^{1/\varepsilon} = p_{r,1}\cdot 0.70^{1/\varepsilon}.
\]
Solving $0.863 = 0.70^{1/\varepsilon}$ yields
\[
\varepsilon = \frac{\ln(0.70)}{\ln(0.863)} \approx 2.4.
\]
Based on our reading of the literature, an elasticity in this range appears quite high for dense urban areas.

Lastly, the initial rent charged to new tenants in tenancy-controlled units $\tilde{p}_{r,0}$, is also set externally to match the observed "entry premium" for new tenants relative to free-market rents. Specifically, we choose $\tilde{p}_{r,0}$ so that new tenants in controlled units pay 5.9\% above free-market rents, consistent with our point estimate for Toronto (95\% confidence interval: 2.7\%–9.1\%), as detailed in Section~\ref{section:appendix_measuring_premium} of the Appendix.

\subsubsection{Mortgages}

In accordance with banking regulations in Canada, we set the LTV constraint, denoted as $\lambda_{1}(ph)$, as follows:

\begin{equation*}
\lambda_{1}(ph) = \begin{cases}
0.8 &\text{if $ph\geq \$1M$}\\
0.9 &\text{if $\$500K\leq ph< \$1M$}\\
0.95 &\text{if $ph < \$500K$}.
\end{cases}
\end{equation*}

\noindent Households purchasing a home valued above a model-equivalent of \$1 million can borrow up to 80\% of its value; those purchasing homes valued between \$500,000 and \$1 million can borrow up to 90\% of the house value; and those purchasing homes valued below \$500,000 face a 95\% LTV cap. We set the PTI constraint parameter $\lambda_{2}$ to $0.39$, consistent with the 39\% PTI threshold commonly used in the Canadian mortgage industry.\footnote{See \url{https://www.cmhc-schl.gc.ca/consumers/home-buying/calculators/debt-service-calculator}.}

Fixed mortgage origination costs $\omega$ are set $0.68\%$ of the benchmark's average household income. This figure is based on the 2016 Census data and the Canada Mortgage and Housing Corporation's (CMHC) website, which includes standard expenses such as appraisal fees ($\approx$\$250) and lawyer fees ($\approx$\$500).\footnote{See \url{https://www.cmhc-schl.gc.ca/en/professionals/industry-innovation-and-leadership/industry-expertise/resources-for-mortgage-professionals/how-much-will-my-home-really-cost}.}

Following the methodology in \citet{kkk-2021}, we set $\upsilon$, which reflects the average annual principal repayment rate under a 30-year amortization schedule, to $1.7\%$. Finally, we set the interest rate premium on mortgage debt $\Delta_{m}$ to 0.19\%. This value corresponds to the historical average spread between 5-year variable mortgage rates and 10-year Government of Canada bond yields over the 2002-2016 period.\footnote{Historical 5-year variable mortgage rates can be accessed at \url{https://www.superbrokers.ca/tools/mortgage-rate-history}, and 10-year bond yields can be accessed at \url{https://www.bankofcanada.ca/rates/interest-rates/canadian-bonds/}.}

\subsection{Internally calibrated parameters}\label{sec:internal_calibration}

After assigning the externally determined parameter values discussed above, we proceed to calibrate another eleven parameters, as outlined in Table \ref{tab:calibration_internal}. These variables include the discount factor, $\beta$; the utility weight on the consumption of housing services, $\gamma$; the strength of the bequest motive, $\mu$; the average income taxation level, $\tau_{\ell}$; the variability in labor income's stochastic component, $\sigma_{e}$; the per-period rent decline for continued controlled tenancy, $\kappa$; the parameter governing how strongly the rent gap affects eviction risk, $\epsilon_1$; the parameter governing the time-use costs for controlled renters, $\varkappa$; the minimum housing unit size for purchase, $\underline{h}_{o}$; the smallest rental size, $\underline{h}_{r}$; and housing services for homeless households, $\underline{h}_{h}$.\footnote{The construction cost scalar, $\varsigma_{1}$, is a free parameter selected to equate housing supply to housing demand in the benchmark equilibrium. This value remains fixed in all counterfactual simulations.}

We collectively calibrate these parameters to match eleven key empirical moments in our Toronto data:

\begin{itemize}
    \item The ratio of aggregate mortgage debt to aggregate house values among homeowners, which stands at 21.5\% according to the 2019 Survey of Financial Security (SFS).
    \item The average rent-to-income ratio, which we estimate at 31.1\% using 2016 Census.
    \item The ratio of aggregate bequests to aggregate net wealth, which is 2.2\%, based on estimates we derive from the 2019 SFS.
    \item The average labor income tax rate, recorded at 19.6\%, from the 2016 Census.
    \item The Gini coefficient of labor income, which we also obtain from the 2016 Census.
    \item The ratio of average controlled rent to free-market rent, which is 0.863 based on estimates by \citet{zhang-gou} for Toronto.
    \item The percentage of renters evicted because landlords can obtain higher rents from new tenants, estimated at 0.1\% based on Toronto data from \citet{LeonIveniuk2020}.
    \item The percentage of renters in the controlled market, which we estimate is 71\% in the 2016 Census.
    \item The homeownership rate, which is 66.5\% in the 2016 Census.
    \item The percentage of renters spending over 30\% of their pre-tax income on rent, standing at 41.4\% according to the 2016 Census.
    \item The percentage of individuals experiencing homelessness, which is documented at 0.26\% according to the ``Homelessness Hub".\footnote{See \url{https://www.homelesshub.ca/community-profile/toronto}.}
\end{itemize}

While all parameters are determined jointly in calibration, each one exerts a more significant influence on specific outcomes. Here is how:
\begin{itemize}
    \item Discount factor, $\beta$: Primarily shapes homeowners' debt levels; lower $\beta$ values lead to higher homeowner debt levels.
    \item Utility weight on housing, $\gamma$: Influences the proportion of renters' income spent on housing; higher $\gamma$ values indicate greater desire for housing services among households, and thus a larger share of income spent on housing.
    \item Bequest motive strength, $\mu$: Affects the ratio of bequests to net wealth; higher $\mu$ values result in relatively more bequests compared to overall net wealth.
    \item Average income taxation level, $\tau_{\ell}$: Determines the average income tax rate; lower $\tau_{\ell}$ values correspond to higher tax burdens on households' incomes.
    \item Variability in labor income's stochastic component, $\sigma_{e}$: Shapes labor income inequality; higher $\sigma_{e}$ values lead to greater dispersion in the labor income distribution and thus higher Gini coefficients.
    \item Annual rent decrease for each additional period a tenancy-controlled tenant remains in the unit, $\kappa$: Governs how much rents decline each period for sitting tenancy-controlled tenants. Higher values of $\kappa$ result in a lower average controlled rent relative to the free-market rent.
    \item Increased risk of eviction due to a larger rent gap, $\epsilon_1$: Determines how strongly the rent gap increases the probability of eviction under tenancy control. Higher values of $\epsilon_1$ lead to a higher eviction rate attributable to landlords seeking higher rents from new tenants.
    \item Time-use cost faced by controlled renters, $\varkappa$: Conditional on our estimated "entry premium" in the data, this affects how many renters live in the free market. Larger values of $\varkappa$ decrease the share of renters who choose to remain in the controlled segment.\footnote{We assume that renters value time at 50\% of their hourly wage. Given this preference, the calibrated rent-control distortion implies that controlled renters incur in monetary terms the equivalent of approximately 10.5 additional minutes per day. This magnitude is consistent with 2016 Census data, which show that, on average, controlled renters in Toronto spend about two minutes more per day commuting than free-market renters.}
    \item Minimum housing unit size for purchase, $\underline{h}_{o}$: Influences the homeownership rate; larger minimum sizes lead to fewer households opting for homeownership.
    \item Minimum rental size, $\underline{h}_{r}$: Governs the proportion of renters spending over 30\% of their pre-tax income on housing; larger minimum sizes result in higher percentages of renters facing high housing cost burdens.
    \item Housing services for homeless households, $\underline{h}_{h}$: Affects the homelessness rate; greater services result in higher rates of homelessness.
\end{itemize}

\subsection{Model performance}\label{sec:non_targeted_statistics}

We assess the performance of our benchmark model by evaluating its ability to replicate aggregate, lifecycle, and distributional moments not targeted in our calibration. 

Panel (a) of Table \ref{tab:aggregate_non_targeted_statistics} reports the model's fit relative to key non-targeted aggregate statistics. The model generates a somewhat higher ratio of the average owned-house value to average household income compared to the data, alongside a somewhat lower percentage of owners with positive mortgage debt. Due to data limitations, we cannot directly estimate the ratio of the average rental unit value to average owner-occupied unit value in Toronto. Hence, we compare our model's moment to a U.S. estimate reported by \citet{chatterjee-eyigungor-2015}, finding a slightly lower ratio in our model. Moreover, the ratio of aggregate net worth to GDP in our model is in line with the same ratio observed in our data, and so are the income differences between renters in the free market and renters in the controlled market and the average size difference between units in each market. Lastly, we note that since approximately 30\% of renters in our benchmark reside in the free market and the utilization rate of the exempt stock is $\phi=46.8\%$, about 13.3\% of the total housing stock in our benchmark is exempt from rent control.

In panel (b), we report the share of renters spending more than 50\% and 75\% of their income on rent. Relative to the data, the model understates the prevalence of cost-burdened renters. This bias, however, works against rent control in our analysis: because decontrol imposes smaller losses on low-income renters in the model than it would if there were more cost-burdened renters, any welfare gains from rent control in our experiments are, if anything, understated.

In panel (c) we show that the model's wealth distribution is quite unequal with a Gini coefficient of 0.47, though lower than the 0.65 observed in our 2019 SFS data. The model understates wealth concentration at the top relative to the data due to the absence of rate-of-return heterogeneity, which \citet{fagereng} identify as a key determinant of upper-tail wealth accumulation. Our model also omits heterogeneous time preferences \citep{ks-1998} and the type of fat-tailed income shocks studied in \citet{mj-shocks}, both of which are known to amplify wealth inequality.

The model generates lifecycle patterns in homeownership rates and LTV ratios that resemble those observed in the data. Panel (a) in Figure \ref{fig:life_cycle} shows similar patterns in homeownership rates, though the model slightly over-predicts homeownership among middle-aged households and under-predicts it for young cohorts. Panel (b) plots LTV ratios, which decline with age in both the data and the model. Young households in the model hold higher LTV ratios than observed in the data, whereas older households exhibit somewhat lower leverage.

Table \ref{tab:tenure_by_income_quintile} reports the shares of households living in the controlled market by income quintile. The second column from the left represents the lowest income quintile, while the first column from the right corresponds to the highest quintile. The first row reports statistics from the 2016 Census, while the second row shows the corresponding estimates in the calibrated model. The model does a decent job relative to the data. However, it modestly understates the share of controlled renters in the bottom income quintile. This happens in the model because some marginally income-constrained households in the bottom quintile who want to avoid homelessness optimally choose to rent in the free market rather than pay the "entry premium" associated with controlled units. As in the case of cost-burdened renters, this feature also plays against rent control in our analysis: because decontrol imposes smaller losses on low-income renters in the model, the welfare gains we estimate from rent control are likely a lower bound.

\section{Experiments}\label{section:experiments}

Using our calibrated model, we conduct four experiments to study how different degrees of redistribution efficiency affect the welfare consequences of rent control. In these experiments, we also explore how rental-market policies affect equilibrium rents, housing prices, rental supply, homeownership, and homelessness. In all of our experiments we compare the different long-run equilibria of the alternative policies to our calibrated benchmark equilibrium. We present the results of our experiments in Table \ref{tab:aggregate_effects}.

In all experiments, the local government budget is balanced by adjusting the property tax rate. Changes in landlord profits are rebated as lump-sum transfers to households in the top 35\% of the wealth distribution, consistent with 2019 SFS data, in which 75\% of households who own secondary properties occupy the top 35\% of the wealth distribution. These transfers are added to taxable income and taxed progressively at the federal level. The resulting excess federal revenues are redistributed in one of three ways: income tax cuts to all households (Experiments 1 and 4), lump-sum transfers to renters and the homeless only (Experiment 2)\footnote{We also conducted an additional experiment in which excess government revenues are redistributed as lump-sum transfers exclusively to renters and homeless households in the bottom two income quintiles. Welfare losses in that experiment are smaller than in Experiment 2 relative to the benchmark, but the qualitative conclusions are unchanged. For brevity, we do not report these results.}, and lump-sum transfers to newborns only (Experiment 3).

Although economists traditionally view rent control as harmful because of the deadweight losses it creates, they may nonetheless accept arguments in its favor that focus on costly transition dynamics out of regulated markets. The concern is that, even if a deregulated steady state ultimately improves welfare, the short-run adjustment can be welfare-reducing \citep{fmn-2023} when rent control is already in effect. Our analysis takes a different angle: we ask whether rent control can improve welfare even in the long-run under various conditions. The results we present below show that, under realistic conditions, both tenancy and unit control can indeed raise welfare even in the long run.

\begin{itemize}

\item \textbf{Experiment 1: Full rent decontrol with income tax cuts vs.\ benchmark.}
\begin{itemize}
    \item \textbf{Question:} Can tenancy control with exemptions generate higher long-run welfare than full decontrol when landlord profits are only partially taxed and excess revenues are redistributed through income tax cuts to the entire local population?
    \item \textbf{Change:} We compare the long-run benchmark economy with tenancy control and partial exemptions to an economy with full rent decontrol.
    \item \textbf{Intuition:} This experiment assesses whether, in the absence of fiscal redistributive efficiency, the redistribution achieved under tenancy control outweighs its distortions to rental supply, housing allocation, eviction risk, and equilibrium prices.
\end{itemize}

\item \textbf{Experiment 2: Full rent decontrol with targeted lump-sum transfers to renters and homeless households vs.\ benchmark.}
\begin{itemize}
    \item \textbf{Question and change:} Same as Experiment~1, but excess government revenues are redistributed as lump-sum transfers exclusively to renters and homeless households.
    \item \textbf{Intuition:} This experiment evaluates how sensitive welfare comparisons are to a more targeted, but still rather imperfect, redistribution scheme.
\end{itemize}

\item \textbf{Experiment 3: Full rent decontrol with lump-sum transfers to newborns only vs.\ benchmark.}
\begin{itemize}
    \item \textbf{Question and change:} Same as Experiment~2, except that excess revenues are redistributed as lump-sum transfers to newborns only.
    \item \textbf{Intuition:} This experiment isolates the role of highly efficient redistribution in shaping long-run welfare outcomes.
\end{itemize}

\item \textbf{Experiment 4: Unit control vs.\ benchmark.}
\begin{itemize}
    \item \textbf{Question:} Can unit control outperform the benchmark tenancy-control regime? If not, can it still outperform decontrol?
    \item \textbf{Change:} We replace our benchmark with tenancy control with unit-based rent control in which we hold fixed the fraction of exempt housing stock, and set the regulated rent charged throughout a tenancy — regardless of tenant turnover — at a level equal to the benchmark’s average controlled rent.
    \item \textbf{Intuition:} The welfare effects of unit control are ambiguous because it trades off different sources of misallocation relative to tenancy control. Unit control delivers large, immediate rent discounts but keeps rents fixed over the tenancy, potentially reducing lock-in and eliminating eviction risk for incumbents. At the same time, controlled units are allocated through lotteries rather than prices, which can increase misallocation at entry.
\end{itemize}

\end{itemize}

\paragraph{Welfare measure} Our welfare metric is the expected ex-ante utility of newborns \citep{ralws-1971,conesa-et-al,gkkoc}. We use this measure to compute welfare gains in non‑housing consumption‑equivalent terms, defined as the fixed proportional increase in non‑housing consumption that would make newborn households indifferent between the counterfactual equilibrium and the benchmark equilibrium. Formally, the non-housing consumption-equivalent welfare effect of a policy change is expressed as:

\begin{equation}
  CE = \left[\frac{\int_{\mathcal{S}}\left(V^{\dagger}_{1}(s)-\omega_{B}\left(q'^{\star}_{J}(s)\right)\right)\mathrm{d}\Psi^{\dagger}_{1}(s)}{\int_{\mathcal{S}}\left(V^{\star}_{1}(s)-\omega_{B}\left(q'^{\star}_{J}(s)\right)\right)\mathrm{d}\Psi^{\star}_{1}(s)}\right]^{\frac{1}{(1-\gamma)(1-\sigma)}}-1,
\end{equation}

\noindent where daggers denote objects in the counterfactual equilibrium, and stars refer to benchmark objects. 

\subsection{Experiment 1: Full rent decontrol with income tax cuts vs.\ benchmark}\label{sec:exp_1}

In our first experiment, we fully decontrol the rental market, redistribute excess revenues through income tax cuts, and assess the resulting effects on welfare and other key outcomes, as shown in the first row of Table \ref{tab:aggregate_effects}. This experiment evaluates whether, in the absence of efficient government redistribution, replacing tenancy-based rent discounts with broad income tax reductions can improve welfare once equilibrium responses are taken into account. All reported changes are relative to the tenancy-control benchmark.

Removing tenancy control would naturally increase landlord profits, resulting in lump-sum transfers equal to 1.81\% of average income to each household in the top 35\% of the wealth distribution. These transfers would be added to taxable income and taxed at each household’s marginal rate, generating enough revenue to finance a 0.85\% reduction in average income taxes. Due to progressivity in the tax schedule, these tax cuts would disproportionately benefit higher-income households.

With the removal of tenure-dependent rent discounts, renters would have weaker incentives to remain in the rental market. As a result, transitions into homeownership would increase, raising the homeownership rate by 1.25 percentage points (p.p.), increasing house prices by 1.67\%, and reducing rental demand at the extensive margin. 

On the other hand, full decontrol would remove the lock-in effect related to accumulated discounts, and thus increase rental consumption among remaining renters as their household size increases over their lifecycle. In the aggregate, the intensive and extensive margins largely offset each other, leaving the equilibrium quantity of rental housing essentially unchanged. Free-market rents nevertheless rise by 1.35\%, reflecting higher house prices, which increase the cost of supplying rental housing (see Equation~\eqref{eq:rental_MC}). Our findings that free-market rents and house prices would rise are consistent with the evidence in \citet{sims-2007} and \citet{app-2014}, who document that following rent decontrol in 1994 in Cambridge, Massachusetts, house prices and rents for previously unregulated units both increased.

Crucially, the increase in free-market rents, together with the removal of the rent control option, would reduce renters’ ability to weather severe income shocks, leading homelessness to rise by 90.17\%. The increase in homelessness would require the local government to raise the property tax rate from 0.44\% to about 0.46\%, thereby further increasing the cost of supplying rental housing.

Overall, newborns’ welfare would decline by 0.65\% under decontrol despite the distortions associated with tenancy control. This decline would reflect worse redistribution toward low-income renters: the income tax cuts financed by higher landlord profits would fail to replicate the redistributive role provided by tenure-based rent discounts. This outcome illustrates a realistic failure of the Second Welfare Theorem, whereby removing a distortion can reduce welfare in the absence of sufficiently targeted redistribution.

\subsection{Experiment 2: Full rent decontrol with targeted lump-sum transfers to renters and homeless households vs.\ benchmark}\label{sec:exp_2}

For our second experiment, we again fully decontrol the rental market, but redistribute excess government revenues as lump-sum transfers exclusively to renters and homeless households. Results are reported in the second row of Table \ref{tab:aggregate_effects}. Relative to Experiment 1, this more targeted redistribution would partially mitigate the adverse welfare effects of decontrol with income tax cuts by directing transfers toward households that are more affected by the removal of tenure-based rent discounts.

Newborns’ welfare would still decline relative to the tenancy-control benchmark, but by less than under income tax cuts. At the same time, the targeted transfers would encourage participation in the rental market. As a result, homeownership would rise less than in Experiment 1, and free-market rents would rise by more, with equilibrium rental quantity supplied rising by about 0.57\% above the benchmark.

Overall, this experiment shows that even under a relatively optimistic, more targeted government redistribution, tenancy control would still raise welfare relative to decontrol.

\subsection{Experiment 3: Full rent decontrol with lump-sum transfers to newborns only vs.\ benchmark}\label{sec:exp_3}

In our next experiment, we again fully decontrol the rental market, but assume that the government can tax away the entirety of the additional landlord profits and redistribute excess revenues as lump-sum transfers exclusively to newborn households. This redistribution scheme is meant to approximate highly efficient targeting by directing resources toward households at the beginning of their life cycle, before housing and tenure decisions are made. We report the results of this experiment in the third row of Table \ref{tab:aggregate_effects}.

Under this policy, newborns’ welfare would increase substantially relative to the tenancy-control benchmark (4.16\%). The lump-sum transfers would relax borrowing constraints early in life, allowing households to transition into homeownership more quickly. As a result, the homeownership rate and house prices would rise by more, while free-market rents would rise by less than in Experiments~1 and~2, and the adverse effects associated with the removal of tenure-based rent protection would be more than offset.

This experiment illustrates that when redistribution is sufficiently well targeted, rent decontrol improves welfare significantly, highlighting the central role of redistribution efficiency in determining the welfare consequences of housing policy.

\subsection{Experiment 4: Unit control vs. benchmark}\label{sec:exp_4}

Lastly, we consider replacing tenancy control with unit control while holding fixed the fraction of exempt housing stock and compare outcomes to the benchmark tenancy-control regime. The regulated rental rate throughout a tenancy under unit control is set equal to the average controlled rent under tenancy control — approximately 13.7\% below the benchmark average free-market rent and 18.52\% below the controlled rent paid by new movers under tenancy control. Under unit control, these rent discounts would be provided immediately upon entry rather than accruing gradually with tenure, shifting the benefits of rent control earlier in the life cycle for lottery winners who opt for a controlled rental. Results are shown in the fourth row of Table \ref{tab:aggregate_effects}.

Because unit control prevents initial rents from adjusting to clear the controlled market, access to controlled units would be rationed through lotteries. We therefore choose $\pi$ — the probability of drawing an opportunity to rent a controlled unit ($\eta=1$) — to clear the controlled market. This randomness implies that allocation at entry is less efficient than under tenancy control. However, because households’ decisions to enter the controlled market and their tenancy duration remain endogenous, and because rent discounts do not grow with tenure under unit control — weakening lock-in effects over a tenancy relative to tenancy control — the allocation of the rental stock might improve over time relative to tenancy control. Our simulations show that the entry-stage misallocation would dominate this dynamic effect: although unit control allocates more controlled units to lower-income households overall (see the third row of Table~\ref{tab:tenure_by_income_quintile}), fewer controlled units would be occupied by the poorest households compared with the tenancy-control benchmark.

By weakening lock-in incentives over the course of a tenancy relative to tenancy control, unit control would encourage renters to adjust their housing consumption more readily over the life cycle — occupying smaller units earlier in life and transitioning to larger units as household needs evolve. Relative to tenancy control, this would increase rental demand along the intensive margin.

By shifting rent discounts earlier in life and weakening incentives to remain in controlled units, unit control would accelerate transitions into homeownership, increasing the homeownership rate by 2.54~p.p. As renters would exit the rental market more quickly, aggregate rental demand would decline along the extensive margin, putting downward pressure on free-market rents. In equilibrium, this extensive-margin effect would dominate the intensive margin, causing free-market rents to fall by 8.01\% relative to the tenancy-control benchmark.

The combination of the potential earlier access to rent discounts and lower equilibrium free-market rents would improve renters’ ability to weather adverse income shocks. As a result, homelessness would decline by 39.91\% relative to the tenancy-control benchmark, and the local government would be able to slightly lower property tax rates.

Overall, the welfare effects of unit control reflect a trade-off between earlier access to rent discounts and increased misallocation due to lottery-based access to controlled units. Combined with the decline in equilibrium free-market rents, these forces would reduce welfare by only 0.05\% relative to the benchmark.

\section{Summary and implications}\label{sec:conclusion}

We develop a dynamic general equilibrium model of rent control with exemptions for new construction, calibrated to Toronto data. Our results show that when excess revenues are redistributed through income tax cuts or lump-sum transfers to renters and the homeless, the redistributive benefits of both tenancy and unit control outweigh their distortionary effects relative to full rent decontrol. In contrast, when excess government revenues are rebated directly to newborns, full decontrol indeed outperforms both the benchmark and unit control.

We also find that while unit control worsens the allocation of units across the income distribution, it reduces free-market rents and improves the allocation of discounts across ages, since even newborn households can benefit from lower rents without waiting for tenure-based discounts to accumulate. As a result, even though unit control creates shortages and tenancy control does not, it reduces homelessness and lowers welfare only slightly relative to the benchmark. 

Two policy implications follow from our findings. First, the desirability of rent control should be evaluated in conjunction with the surrounding fiscal environment rather than in isolation. In practice, we think that the most realistic benchmark is the one closest to the income-tax-cut experiment, since federal revenues are typically redistributed broadly across households, including those living outside the regulated locality. Second, if the government is willing to accept a modest welfare loss relative to tenancy control, unit control seems to be a better tool for reducing homelessness more substantially than tenancy control. 

\bibliographystyle{aer}
\bibliography{reference_files/bib}

\vspace{1cm}

\begin{table}[H]
    \centering
    \caption{Externally assigned parameter values}
    \vspace{-0.5cm}
    \label{tab:calibration_external}
    \input{tables/calibration_external}
    \vspace{-0.45cm}
    \begin{fignote2}{1.02\textwidth}
    \textit{Notes:} Table \ref{tab:calibration_external} presents the values of the externally calibrated parameters. Part (i) delineates the external calibration for non-housing parameters, while part (ii) outlines the external calibration of housing parameters. Part (iii) provides the external calibration of mortgage-related parameters. $\overline{y}_{\ell}$ denotes the average labor income and $\overline{h}_{c}$ denotes the benchmark equilibrium's average free-market rental size.
    \end{fignote2}
\end{table}

\begin{table}[H]
    \centering
    \caption{Internally assigned parameter values}
    \vspace{-0.5cm}
    \label{tab:calibration_internal}
    \input{tables/calibration_internal}
    \vspace{-0.45cm}
    \begin{fignote2}{0.87\textwidth}
    \textit{Notes:} Table \ref{tab:calibration_internal} displays the values of the internally calibrated parameters. Our benchmark is closely aligned with all of our targets. From the 2019 SFS, we compute the percentage of homeowners with mortgages and the aggregate bequests to aggregate net wealth ratio. From the 2016 Census, we obtain data for Toronto CMA regarding the total rent to total income ratio, the average income tax rate, the labor income Gini coefficient, the percentage of renters in the controlled market, the percentage of renters spending over 30\% of their income on rent, and the homeownership rate. The average controlled rent to free-market rent comes from \citet{zhang-gou}, and the eviction rate due to tenancy control is estimated using data from \citet{LeonIveniuk2020}. Lastly, we obtain the homelessness rate from the ``Homelessness Hub". $\overline{h}_{o}$ denotes the average owned-house size. $\overline{y}_{\ell}$ denotes the average labor income.
    \end{fignote2}
\end{table}

\begin{table}[H]
    \centering
    \caption{Aggregate non-targeted statistics}
    \vspace{-0.5cm}
    \label{tab:aggregate_non_targeted_statistics}  \input{tables/aggregate_non_targeted_statistics.tex}
    \vspace{-0.45cm}
    \begin{fignote2}{0.84\textwidth}
      \textit{Notes:} Table \ref{tab:aggregate_non_targeted_statistics} presents the performance on non-targeted aggregate statistics compared to the data. The second column from the left displays the benchmark model's performance, the third column includes estimates from the literature or the data, and the last column provides the source.
    \end{fignote2}
\end{table}

\begin{table}[H]
    \centering
    \caption{Households in controlled market by income quintile (all households)}
    \vspace{-0.5cm}
    \label{tab:tenure_by_income_quintile}     \input{tables/tenure_by_income_quintile_everyone.tex}
    \vspace{-0.45cm}
    \begin{fignote2}{0.81\textwidth}
    \textit{Notes:} Table \ref{tab:tenure_by_income_quintile} displays the percentage of households renting in the controlled market by income quintile. The first row presents these statistics in the data, which we computed using data from the 2016 Census. The second row displays outcomes in our model's benchmark with tenancy control, and the last row represents the performance of our model with unit control instead of tenancy control. The second column from the left represents the lowest income quintile, while the first column from the right signifies the highest quintile.  
    \end{fignote2}
\end{table}

\begin{table}[H]
    \centering
    \caption{Aggregate effects generated by changes in rent-control policy}
    \vspace{-0.5cm}
    \label{tab:aggregate_effects}
    \input{tables/aggregate_effects}
    \vspace{-0.45cm}
    \begin{fignote2}{0.85\textwidth}
    \textit{Notes:} Table \ref{tab:aggregate_effects} compares the benchmark with a series of long-run equilibria under alternative rental market policies: full rent decontrol with income tax cuts, decontrol with targeted lump-sum transfers, decontrol with transfers to newborns only, and a unit-control regime in which the exempt housing share is the same as in the benchmark and the regulated rental rate is equal to the benchmark average controlled rent.
    \end{fignote2}
\end{table}

\begin{figure}[H]
  \centering
    \caption{Diagram illustrating rent trajectories under tenancy control based on rental market selection and duration in the same rental unit}
    \vspace{-0.37cm}
    \label{fig:diagram_rent_path}
    \includegraphics[width=\textwidth]{figures/path_of_rent_diagram.pdf}
    \vspace{-1.3cm}
    \begin{fignote2}{1.0\textwidth}
    \textit{Notes:} The downward-sloping blue line shows how rents decline over time for tenants in tenancy-controlled units, starting from the initial rent $\tilde{p}_{r,0}$ and averaging to $\overline{\tilde{p}}_{r,0}$. The horizontal red line represents the free-market rent, $p_{r,1}$. Free-market rents remain constant regardless of tenancy duration. Initially, controlled tenants pay more, but over time they benefit from lower rents. This "entry premium" deters short-term tenants — who prefer lower upfront costs — from choosing controlled units, while long-term tenants are willing to pay it for improved affordability later.
   \end{fignote2}
\end{figure}

\begin{figure}[H]
  \centering
    \caption{Lifecycle statistics}
    \vspace{-0.2cm}
    \label{fig:life_cycle}
    \includegraphics[width=\textwidth]{figures/homeownership_and_LTV.png}
    \vspace{-0.45cm}
    \begin{fignote2}{1\textwidth}
      \textit{Notes:} In Figure \ref{fig:life_cycle}, the two panels illustrate the model's performance on lifecycle statistics compared to the data. Panel (a) displays the average homeownership rate for each age cohort, while panel (b) presents the average loan-to-value (LTV) ratios for homeowners across 12 age groups, as the 2019 SFS provides household maintainers' ages in 5-year ranges. Model outputs are depicted in orange, while data are represented in blue. We source homeownership rates by age for Toronto CMA from the 2016 Census, and obtain data on the average LTV ratios from the 2019 SFS.
    \end{fignote2}
\end{figure}

\newpage

\include{appendix}

\end{document}